\documentclass[11pt]{article}

\usepackage[margin=1.2in]{geometry}
\usepackage{amsmath,amssymb,amsthm}
\usepackage{booktabs}
\usepackage{enumitem}
\usepackage[colorlinks=true,linkcolor=blue,citecolor=blue,urlcolor=blue]{hyperref}

\newtheorem{theorem}{Theorem}[section]
\newtheorem{lemma}[theorem]{Lemma}
\newtheorem{proposition}[theorem]{Proposition}
\newtheorem{corollary}[theorem]{Corollary}
\newtheorem{conjecture}[theorem]{Conjecture}
\theoremstyle{definition}
\newtheorem{definition}[theorem]{Definition}
\newtheorem{example}[theorem]{Example}
\newtheorem{question}[theorem]{Question}
\theoremstyle{remark}
\newtheorem{remark}[theorem]{Remark}
\newtheorem{observation}[theorem]{Observation}

\newcommand{\mass}{\operatorname{mass}}
\newcommand{\ord}{\operatorname{ord}}
\newcommand{\maxord}{\operatorname{maxord}}
\newcommand{\N}{\mathbb{N}}
\newcommand{\intv}[1]{[#1]}
\newcommand{\Otilde}{\widetilde{O}}

\title{Forbidden sequences in the sumset partition problem}
\author{Angel Pelayo\thanks{Email: \texttt{angel.pelayo@umich.mx}. Code reproducing all
computational claims is available at
\url{https://github.com/AngelMisaelPelayo/sumset-partition-forbidden-sequences}.}}
\date{July 20, 2026}

\begin{document}
\maketitle

\begin{abstract}
A non-increasing sequence $M=(m_1\ge\dots\ge m_k)$ of positive integers is
\emph{realizable} if there are pairwise disjoint subsets
$X_1,\dots,X_k\subseteq\{1,\dots,m_1\}$ with $\sum_{x\in X_i}x=m_i$, and
\emph{forbidden} otherwise; this relaxes the classical sumset partition problem,
where the sets must partition $\{1,\dots,n\}$ exactly. Realizability is
hereditary under subsequences, so forbidden sequences are generated by
\emph{minimal} ones. A universal singleton rule shows that sequences with
distinct entries are always realizable and reduces the decision problem to a
packing problem for repeated values --- conjecturally, to a full rainbow matching
problem over arithmetically structured color classes. We catalog all $2{,}668$
minimal forbidden sequences with $m_1\le12$, exhibit an infinite extremal
family, and prove that no minimal forbidden sequence lies on the
exact-partition boundary. Returning to exact $n$-realizability, tree search
pruned by the relaxed oracle is a complete decision procedure, and its
refutation depth defines an \emph{order} on feasible sequences that are not
$n$-realizable.
The resulting hierarchy is unbounded:
$\lfloor(n-1)/2\rfloor\le\maxord(n)\le n-5$, with equality on the right for all
$8\le n\le31$, and two horizon theorems show that no witness family with a
fixed finite core can certify the equality for all $n$. The computational
complexity of both realizability notions remains open; we isolate a structured
rainbow matching problem whose resolution would settle the relaxed case.
\end{abstract}

\section{Introduction}\label{sec:intro}

Let $\intv{n}=\{1,\dots,n\}$. The \emph{sumset partition problem} asks which
prescribed block sums are achievable by partitions of $\intv{n}$.

\begin{definition}\label{def:nrealizable}
A non-increasing sequence $M=(m_1\ge m_2\ge\dots\ge m_k)$ of positive integers is
\emph{$n$-realizable} if $\intv{n}$ admits a partition $X_1\cup\dots\cup X_k$ with
$\sum_{x\in X_i}x=m_i$ for each $i$. The sequence $M$ is \emph{$n$-feasible} if
$\sum_i m_i=\binom{n+1}{2}$.
\end{definition}

Feasibility is an obvious necessary condition for realizability, and the problem
is to determine which $n$-feasible sequences are $n$-realizable.

The question originates in the Ascending Subgraph Decomposition (ASD) conjecture of
Alavi, Boals, Chartrand, Erd\H{o}s and Oellermann~\cite{AlaviBCEO87}: $n$-realizability
of $M$ is precisely the existence of a star-forest ASD of a star forest with star
sizes $m_1,\dots,m_k$. The positive literature is a sequence of successively weaker
sufficient conditions: Ma, Zhou and Zhou~\cite{MaZhouZhou94} proved that $m_k\ge n$
suffices, confirming the star-forest case of the ASD conjecture; Chen, Fu, Wang and
Zhou~\cite{ChenFuWangZhou05} weakened this to $m_{k-1}\ge n$; Lladó and
Moragas~\cite{LladoMoragas09} characterized $n$-realizability for $k\le 4$ and proved
that for $k\ge5$ the conditions $n\ge 4k-1$ and $m_k>4k-1$ together suffice, with a
modular analogue via the polynomial
method in~\cite{LladoMoragas12}; Aroca and Lladó~\cite{ArocaLlado18} extended the
star-forest theory to bipartite ASD. The ASD conjecture itself was recently proved for
all sufficiently large $n$ by Katsamaktsis, Letzter, Pokrovskiy and
Sudakov~\cite{KLPS25}. Related prescribed-sum partition problems appear in graph
burning~\cite{TanTeh24}, distance-magic labelings~\cite{CichaczGorlich18}, and the
equal-sums ``cutting sticks'' literature~\cite{BuchelGW18,HooryKotlar26}.

All of these results live on the \emph{sufficiency} side: they identify regimes where
no obstruction exists. The present paper develops the \emph{obstruction} side. Our
starting point is a relaxation that turns out to carry all of the small-element
structure of the problem while discarding the rigid covering constraint:

\begin{definition}\label{def:realizable}
$M=(m_1\ge\dots\ge m_k)$ is \emph{realizable} if there are pairwise disjoint subsets
$X_1,\dots,X_k\subseteq \intv{m_1}$ with $\sum_{x\in X_i}x=m_i$ for all $i$; otherwise
$M$ is \emph{forbidden}. A subsequence $S=(m_i)_{i\in I}$ is judged intrinsically,
against its own ground set $\intv{\max S}$.
\end{definition}

The notions of \emph{forbidden} and (\emph{simply}) \emph{realizable} sequences
originate with Lladó and Moragas: they were introduced, together with the minimal
forbidden sequences of length at most $3$, in the announcement~\cite{LladoMoragas09}
and Moragas's thesis~\cite[\S4.2.1]{Moragas10}, and the modular
paper~\cite{LladoMoragas12} proved that for $n\ge(4k)^3$ an $n$-feasible sequence
of length $k$ is $n$-realizable if and only if it contains no forbidden subsequence of
elements smaller than $n$ --- forbidden subsequences being the natural obstruction to
$n$-realizability. It is precisely this characterization that motivates studying
simple realizability in its own right and connecting it back to $n$-realizability.

Note that no covering and no feasibility equation is imposed; only the volume
inequality $\sum_i m_i\le\binom{m_1+1}{2}$ survives as a necessary condition. The
constant-sum case of Definition~\ref{def:realizable} was solved by Ando, Gervacio and
Kano~\cite{AndoGervacioKano90}: $\intv{n}$ contains $k$ disjoint subsets of equal sum
$m$ if and only if $2k-1\le m\le n(n+1)/(2k)$ (see also~\cite{FuHu94,EnomotoKano95}
for the odd-ground-set variants). Beyond its role as an obstruction to
$n$-realizability in~\cite{LladoMoragas09,Moragas10,LladoMoragas12} --- points of
contact we detail in Sections~\ref{sec:singleton} and~\ref{sec:minimal} --- the
general prescribed-sums relaxation appears not to have been studied in its own
right.

\subsection{Results}

\textbf{Heredity and collapse (Section~\ref{sec:heredity}).}
Realizability is downward closed under subsequences
(Corollary~\ref{cor:heredity}). Consequently the natural ``pattern containment''
question collapses: $M$ contains a forbidden subsequence if and only if $M$ itself is
forbidden (Theorem~\ref{thm:collapse}), and the forbidden tails of $M$ form an initial
segment. Forbidden sequences thus form an up-set in the subsequence order, generated
by its minimal elements.

\textbf{The singleton rule and the core (Section~\ref{sec:singleton}).}
A universal exchange argument (Theorem~\ref{thm:singleton}) shows that committing the
singleton $\{v\}$ to an entry of value $v$ never destroys realizability. Two
consequences: sequences with pairwise distinct entries are always realizable ---
forbiddenness is purely a phenomenon of repeated values --- and the decision problem
reduces to a \emph{core}: placing the excess copies of repeated values as disjoint
subsets, of size at least two, of the ground set minus the distinct values. We derive
Hall-type counting conditions (Section~\ref{subsec:hall}), show that greedy rules
provably fail (Section~\ref{subsec:greedy}), and reduce the core, under a
pairs-suffice conjecture verified for all sequences with $m_1\le18$
(Conjecture~\ref{conj:pairs}), to a full rainbow matching problem whose color classes
are the arithmetically structured menus
$\{\{v\}\}\cup\{\{x,v-x\}:1\le x<v/2\}$ (Section~\ref{subsec:rainbow}). General full
rainbow matching is NP-hard even when every color class is a matching of size
two~\cite{ItaiRodehTanimoto78,LePfender14,HommelsheimJM26}, so any tractability must
come from this arithmetic structure --- a sharp open question isolated by the
reduction.

\textbf{Minimal forbidden sequences (Section~\ref{sec:minimal}).}
Every minimal forbidden sequence begins with a doubled maximum, and more strongly,
realizability depends only on the entries at most the largest repeated value
(First-Repeat Reduction, Theorem~\ref{thm:firstrepeat}). We enumerate all minimal
forbidden sequences with $m_1\le 12$ --- there are $2{,}668$, with counts by maximum
entry $1,1,3,3,11,12,42,57,158,291,734,1355$, a sequence not in the
OEIS~\cite{OEIS} as of this writing --- and prove that the family
$F_j=(4j,4j,\underbrace{2j-1,\dots,2j-1}_{j})$ is minimal forbidden for every $j\ge1$
(Theorem~\ref{thm:extremal}), giving length-$k$ minimal forbidden sequences with
maximum entry $4(k-2)$ --- one more than the $4k-9$ of
Moragas~\cite{Moragas10}. We conjecture this is extremal
(Conjecture~\ref{conj:frontier}). Finally, no minimal forbidden sequence lies on the
exact-partition boundary $\sum_im_i=\binom{m_1+1}{2}$
(Theorem~\ref{thm:boundary}): on the boundary, forbiddenness always comes from a
proper subsequence with positive slack.

\textbf{The assignment tree and the order hierarchy
(Sections~\ref{sec:tree}--\ref{sec:order}).}
We then return to exact $n$-realizability. The \emph{assignment tree} greedily assigns
the largest unplaced element $n$ to an entry and recurses; pruning at nodes that are
forbidden \emph{in the relaxed sense} is sound, and we prove the resulting procedure
is complete: $M$ is $n$-realizable iff some maximal path consists entirely of
(relaxed-)realizable nodes (Theorem~\ref{thm:tree}). For $n$-feasible sequences that
are \emph{not} $n$-realizable this yields a natural invariant, the \emph{order}: the
maximum over maximal paths of the level of the first forbidden node --- equivalently,
the refutation depth of the relaxed oracle. Order $1$ means the relaxation already
detects non-realizability at the root; order $\ge2$ sequences are realizable in the
relaxed sense but not $n$-realizable, and we show they must satisfy $m_1\ge n+2$ and
$m_2\ge n+1$ (Theorems~\ref{thm:swap} and~\ref{thm:secondentry}). Both bounds are
attained simultaneously at every level by an explicit order-$2$ family $D_n$ with
pairwise distinct entries (Proposition~\ref{prop:Dn}), the case $n=6$ being
$(8,7,3,2,1)$. Distinct-entry witnesses --- automatically of order $\ge2$ --- are
stratified further by a Hall-type counting oracle for the exact problem
(Proposition~\ref{prop:counting}), incomparable with the relaxed oracle: $D_n$ is
refuted by counting alone, whereas the all-odd \emph{odd-towers} family $O_n$
satisfies every counting condition yet fails $n$-realizability through a parity
mechanism, with $\lfloor n/4\rfloor\le\ord(O_n)\le n-5$ and computed order $n-8$
for $24\le n\le32$ (Theorem~\ref{thm:oddtowers}).

Writing $\maxord(n)$ for the largest order at level $n$, our main results on the
hierarchy are:
\begin{itemize}[itemsep=1pt]
\item $\maxord(n)\le n-5$ for all $n\ge8$ (Theorem~\ref{thm:upper}), by a two-line
induction from an exhaustive base case;
\item an explicit two-parameter ``floor-and-towers'' family $W_a$ with
$\ord(W_a)=a$ exactly, at $n=2a+1$ (Theorem~\ref{thm:Wa}); hence every positive
integer is the order of some sequence and
$\maxord(n)\ge\lfloor(n-1)/2\rfloor$ (Corollary~\ref{cor:unbounded});
\item $\maxord(n)=n-5$ exactly for $8\le n\le31$: the witness sets are enumerated
exhaustively for $n\le23$ by an exact lifting scan (Proposition~\ref{prop:lift},
Table~\ref{tab:lift}), and a capped scan certifies individual witnesses --- hence the
equality, though no longer the exact counts --- for $24\le n\le31$;
\item two \emph{horizon theorems} (Theorems~\ref{thm:topsingleton}
and~\ref{thm:horizon}) showing that every witness family built from a fixed finite
core becomes relaxed-forbidden beyond a computable level --- so a proof of
$\maxord(n)=n-5$ in general requires witnesses whose structure scales with $n$.
\end{itemize}

\textbf{Complexity (Section~\ref{sec:complexity}).}
The complexity of deciding $n$-realizability --- and of the relaxed problem --- appears
to be open; no NP-hardness proof or polynomial algorithm is known even with the input
in unary. The standard strong NP-hardness sources (3-Partition~\cite{GareyJohnson79},
multiple subset sums~\cite{CapraraKP00,CieliebakEPS08}) all use engineered multisets
of items, incompatible with the ground set being exactly $\intv{n}$; for each fixed
$k$ the problem is polynomial by $k$-Subset-Sum dynamic
programming~\cite{AntonopoulosPPV23}. For the relaxed problem we prove two
parameterized upper bounds: unconditionally, realizability is decidable in
randomized time $\Otilde(k+v_1^{\,r})$, where $r=k-d$ is the number of excess
copies of repeated values and $v_1$ the largest repeated value
(Proposition~\ref{prop:coredp}) --- the singleton rule collapses the exponent from
$k$ to $r$, a collapse the covering constraint of $n$-realizability forbids ---
and, under the pairs conjecture, in randomized time
$\mathrm{poly}(m_1)\cdot(k+1)^{O(d)}$ via Exact Matching
(Proposition~\ref{prop:exactmatching}), polynomial for any bounded number $d$ of
distinct values. We collect what is known, describe an exact
solver that classifies all sequences with $m_1\le12$ in seconds, and state the
structured rainbow matching question that our reductions isolate.

\subsection{Relation to prior work}

Closest to our obstruction-side viewpoint is the work of Lladó and Moragas
already discussed~\cite{LladoMoragas09,Moragas10,LladoMoragas12}. Beyond
introducing forbidden sequences, Moragas's thesis~\cite[\S4.2.1]{Moragas10}
determines the minimal forbidden sequences of length at most $3$, bounds the
largest entry of a length-$k$ minimal forbidden sequence from below by $4k-9$,
and conjectures the upper bound $4k$; our Theorem~\ref{thm:extremal} and
Conjecture~\ref{conj:frontier} sharpen both statements, and our pairs conjecture
(Conjecture~\ref{conj:pairs}) strengthens Conjecture~4.12 of~\cite{Moragas10}
(see Sections~\ref{subsec:rainbow} and~\ref{sec:minimal}).
Their obstructions live in the exact-partition problem for fixed $k$ and very large
$n$; our forbidden sequences are intrinsic to the relaxed problem, where heredity
makes the obstruction theory hereditary and (Theorem~\ref{thm:collapse}) containment
collapses to membership. In the graph-burning setting, Tan and Teh~\cite{TanTeh24}
study prescribed-sum partitions of the first $m$ odd integers and conjecture that
beyond an explicit threshold all non-realizable instances are ``trivially''
non-realizable --- a phenomenon parallel to our volume conditions and forbidden
catalog. On the pairs side, Skolem sequences~\cite{Skolem57} and Langford
pairings~\cite{Langford58,Davies59} concern partitions of $\intv{2n}$ into pairs
with prescribed \emph{differences}; the prescribed-\emph{sums} analogue implicit in
Conjecture~\ref{conj:pairs} appears to be new.

\subsection{Notation and the role of computation}

Throughout, $M=(m_1\ge\dots\ge m_k)$ is non-increasing with positive integer entries;
$T_i=(m_i,m_{i+1},\dots,m_k)$ denotes the $i$-th tail. We write
$\mass(S)=\sum_{x\in S}x$ for finite $S\subseteq\N$, so
$\mass(\intv{n})=\binom{n+1}{2}$.

Computer search enters this paper in three roles, which we keep separate.
First, exhaustive enumeration whose output is itself the result: the catalog of
minimal forbidden sequences (Table~\ref{tab:catalog}) and the census of the
order hierarchy (Appendix~\ref{app:census}). Second, evidence for the
conjectures we state (Conjectures~\ref{conj:pairs}, \ref{conj:frontier}
and~\ref{conj:lower}), always reported with its exact search range. Third, two
results consume finite computations as ingredients of their proofs: the
exhaustive base case $\maxord(8)=3$ of Theorem~\ref{thm:upper}, and the
per-witness certificates extending $\maxord(n)=n-5$ from $n=23$ to $n=31$
after Proposition~\ref{prop:lift}. Every other proof is machine-independent;
computation served only to discover the examples and families whose properties
are then proved. All searches are reproducible from the scripts in the
repository cited on the title page, and each computational claim names the
script producing it.

\section{Heredity and the collapse of subsequence containment}\label{sec:heredity}

The relaxed problem owes its structure to one triviality.

\begin{lemma}[Self-bounding supports]\label{lem:selfbound}
In any realization of $M$, every element of $X_i$ is at most $m_i$.
\end{lemma}

\begin{proof}
$X_i$ is a set of distinct positive integers summing to $m_i$; each element is at
most the total.
\end{proof}

\begin{corollary}[Heredity]\label{cor:heredity}
If $M$ is realizable and $S=(m_i)_{i\in I}$ is any subsequence of $M$, then $S$ is
realizable.
\end{corollary}

\begin{proof}
Keep the sets $X_i$, $i\in I$. They are pairwise disjoint and realize the entries of
$S$; by Lemma~\ref{lem:selfbound} each lies in $\intv{m_{i_1}}$ where $i_1=\min I$,
and $m_{i_1}=\max S$ since $M$ is sorted. So the kept sets live in the intrinsic
ground set of $S$.
\end{proof}

Heredity is where sortedness and the intrinsic ground-set convention do their work:
deleting the largest entry shrinks the ground set, and without
Lemma~\ref{lem:selfbound} nothing would guarantee the surviving sets fit.

Realizability being downward closed, forbidden sequences form an up-set in the
subsequence order. This immediately collapses the containment question that motivated
this work:

\begin{theorem}[Collapse]\label{thm:collapse}
For every non-increasing $M$ the following are equivalent:
\begin{enumerate}[label=(\alph*),itemsep=1pt]
\item $M$ contains a forbidden subsequence;
\item some tail $T_i$ is forbidden;
\item $M$ itself is forbidden.
\end{enumerate}
Moreover, if $S=(m_i)_{i\in I}$ is forbidden then already the tail $T_{\min I}$ is
forbidden, and $\{i: T_i \text{ forbidden}\}$ is an initial segment of
$\{1,\dots,k\}$.
\end{theorem}

\begin{proof}
(c)$\Rightarrow$(a): $M$ is a subsequence of itself.
(a)$\Rightarrow$(b): let $S$ be forbidden with first index $i_1$. Then $S$ is a
subsequence of $T_{i_1}$ with the same maximum entry, hence the same ground set; if
$T_{i_1}$ were realizable, Corollary~\ref{cor:heredity} would realize $S$. So
$T_{i_1}$ is forbidden.
(b)$\Rightarrow$(c): $T_i$ is a subsequence of $M$; apply
Corollary~\ref{cor:heredity} in contrapositive.
For the initial segment: if $i\le j$ then $T_j$ is a subsequence of $T_i$, so $T_j$
forbidden implies $T_i$ forbidden.
\end{proof}

Thus ``does $M$ contain a forbidden subsequence?'' is literally the question ``is
$M$ forbidden?'' --- no subsequence search, linear or exponential, is ever needed. The
entire computational content of the relaxed theory is the single decision problem of
Definition~\ref{def:realizable}, taken up in Sections~\ref{sec:singleton}
and~\ref{sec:complexity}. By Theorem~\ref{thm:collapse} the threshold index
$r=\max\{i:T_i\text{ forbidden}\}$ is well defined; when $M$ is forbidden, the longest
realizable tail can be found with $O(\log k)$ oracle calls by binary search.

A first family of necessary conditions also falls out of
Lemma~\ref{lem:selfbound}: all sets realizing $T_i$ fit inside $\intv{m_i}$, so
realizability of $M$ requires the \emph{tail volume conditions}
\begin{equation}\label{eq:volume}
\sum_{j\ge i} m_j \;\le\; \binom{m_i+1}{2}\qquad\text{for every } i.
\end{equation}
These are far from sufficient: $M=(6,6,2,1)$ satisfies \eqref{eq:volume}, with
equality at $i=3$ and $i=4$, yet is forbidden: in $\intv6$ the entries $1$ and $2$ force
$X=\{1\}$ and $X=\{2\}$, after which only one subset summing to $6$ avoids
$\{1,2\}$, so two $6$'s cannot coexist. Of the $130$ minimal forbidden sequences
with $m_1\le8$ (Section~\ref{sec:minimal}), $63$ pass all volume conditions:
\emph{which} small elements get consumed matters beyond volume.

\section{The singleton rule, the core, and pairs}\label{sec:singleton}

By Theorem~\ref{thm:collapse} the relaxed theory rests on a single decision
problem, and the example $(6,6,2,1)$ has already exhibited its mechanism: entries
compete for small elements through forced commitments. This section turns the
mechanism into structure. An exchange rule shows that singleton commitments are
always free, reducing the problem to a core packing problem for repeated values;
counting conditions and greedy rules then probe the core and fail in instructive
ways; and a pairs conjecture, verified far beyond the range of our catalogs,
would reduce the core to a structured rainbow matching problem.

\subsection{The universal singleton rule}

A \emph{partial realization} of $M$ is a family $\mathcal{P}=(X_i)_{i\in J}$,
$J\subseteq\{1,\dots,k\}$, of pairwise disjoint subsets of $\intv{m_1}$ with
$\mass(X_i)=m_i$ for $i\in J$; its blocks, and the elements they contain, are
\emph{committed}. Call $\mathcal{P}$ \emph{extendable} if it extends to a full
realization of $M$ (so $M$ is realizable iff the empty partial realization is
extendable).

\begin{theorem}[Singleton rule]\label{thm:singleton}
Let $\mathcal{P}$ be an extendable partial realization of $M$, let $j\notin J$ be an
uncommitted entry of value $v=m_j$, and suppose the element $v$ is uncommitted. Then
$\mathcal{P}$ together with $X_j=\{v\}$ is again extendable. In particular, if $M$ is
realizable, then it has a realization assigning the singleton $\{v\}$ to one copy of
$v$.
\end{theorem}

\begin{proof}
Take a realization $(X_i)_{i=1}^{k}$ extending $\mathcal{P}$. If the element $v$ is
unused, replace $X_j$ by $\{v\}$. Otherwise $v\in X_{j'}$ for some $j'$, and
$j'\notin J$ since $v$ is uncommitted. If $X_{j'}=\{v\}$ then $m_{j'}=v$; swap the
blocks of the (uncommitted) entries $j$ and $j'$. Otherwise replace $X_{j'}$ by
$(X_{j'}\setminus\{v\})\cup X_j$: the sum is unchanged ($m_{j'}-v+v=m_{j'}$), all
elements are at most $\max(m_{j'},v)=m_{j'}$ (as $v\in X_{j'}$ gives $v\le m_{j'}$ by
Lemma~\ref{lem:selfbound}), and disjointness is preserved since
$X_j\cap X_{j'}=\emptyset$; then set $X_j=\{v\}$. In every case only uncommitted
blocks were modified, so the result still extends $\mathcal{P}$.
\end{proof}

\begin{corollary}\label{cor:distinct}
A sequence with pairwise distinct entries is realizable (all singletons).
Forbiddenness is purely a phenomenon of repeated values.
\end{corollary}

\begin{corollary}[Core reduction]\label{cor:core}
Let $V$ be the set of distinct values of $M$, and for $v\in V$ let $\mu(v)$ be its
multiplicity. Then $M$ is realizable iff the \emph{core} problem is solvable: place,
for each $v\in V$, $\mu(v)-1$ pairwise disjoint subsets of
$A=\intv{m_1}\setminus V$, each of size $\ge2$ and sum $v$.
\end{corollary}

\begin{proof}
Starting from the empty partial realization, iterate Theorem~\ref{thm:singleton} over
one copy of each distinct value: at each step the committed elements are exactly the
values already handled, so the element $v$ is uncommitted at its turn. After the
iteration the committed elements are exactly $V$. An excess copy of $v$ cannot use
the element $v$ (committed), so its set has size $\ge2$.
\end{proof}

Every excess copy of $v$ needs two elements of $A$ summing to $v$ or more elements
yet, hence consumes at least one element $\le\lfloor(v-1)/2\rfloor$: realizability is
governed by how excess copies compete for the \emph{small} numbers of $A$. The
covering-free relaxation has discarded everything else; compare
Section~\ref{sec:tree}, where the discarded large-number placement resurfaces as the
order hierarchy.

\subsection{Hall-type necessary conditions}\label{subsec:hall}

Corollary~\ref{cor:core} yields cheap counting bounds. For every threshold $t\ge1$:
\begin{align}
\text{(count)}\quad
&2\cdot\#\{\text{excess copies with } v\le t\}
+\#\{\text{excess copies with } v>t,\ \lfloor(v-1)/2\rfloor\le t\}
\;\le\; |A\cap\intv{t}|,\notag\\
\text{(mass)}\quad
&\sum\{v : \text{excess copies with } v\le t\} \;\le\; \mass(A\cap\intv{t}).\notag
\end{align}
Empirically (all $m_1\le10$; $309{,}818$ sequences): the volume
conditions~\eqref{eq:volume} alone reject $301{,}848$; of the $3{,}493$ forbidden
sequences surviving volume, the Hall pair catches $2{,}157$ ($62\%$) with no false
rejections. The remaining $1{,}336$ escape for structural reasons --- e.g.\
$(5,5,3,1)$, whose core demands one $5$ from $A=\{2,4\}$: the count of small
elements is right, but the partner of $2$ (namely $3$) is absent. Which small
numbers are present matters, not just how many.

\subsection{Greedy fails}\label{subsec:greedy}

No simple priority rule decides the core. Processing demands in decreasing order,
pairing $v$ with the smallest available $x$ (partner $v-x$), errs on $178$ of the
$4{,}477$ realizable sequences with $m_1\le10$; the ``largest $x<v/2$'' variant errs
on $151$. A clean witness is $M=(10,10,9,9,6,5,4)$, whose core demands $\{10,9\}$ on
$A=\{1,2,3,7,8\}$: the unique solution is $10=\{3,7\}$, $9=\{1,8\}$, and the greedy
pair $\{2,8\}$ for $10$ steals $8$, the only usable partner of $9$. Choices for
different values are genuinely coupled.

\subsection{Pairs suffice (conjecturally), and rainbow matchings}\label{subsec:rainbow}

The coupled choices behind these failures always involve \emph{pairs}
$\{x,v-x\}$, and empirically, pairs are all a realization ever needs:

\begin{conjecture}[Pairs]\label{conj:pairs}
Every realizable sequence admits a realization with $|X_i|\le2$ for all $i$.
\end{conjecture}

Conjecture~\ref{conj:pairs} has been verified for \emph{all} $7{,}423{,}476$
non-increasing sequences with $m_1\le12$: every volume-feasible sequence was solved
exactly, and each of the $30{,}374$ realizable ones was re-solved under the
size-$\le2$ restriction, with no exceptions. Two caveats make it non-obvious. First,
the exchange in Theorem~\ref{thm:singleton} can \emph{grow} another set by one
element, so stripping is not sound under a size cap: the conjecture concerns
unstripped realizations. Second, after stripping, triples can be forced: the core of
$(6,6,5,4)$ demands a $6$ from $A=\{1,2,3\}$, realizable only as $\{1,2,3\}$ --- yet
the unstripped instance has the pairs realization $\{6\},\{2,4\},\{5\},\{1,3\}$.
Since a realization with all $|X_i|\le2$ uses at most $2k$ elements,
Conjecture~\ref{conj:pairs} strengthens Conjecture~4.12 of
Moragas~\cite{Moragas10} --- every realizable sequence admits a realization with
$|\bigcup_iX_i|\le2k$ --- which is shown there to imply the frontier bound
$n_{\max}(k)<4k$ of Section~\ref{sec:minimal}.

Under Conjecture~\ref{conj:pairs}, realizability is exactly a \emph{full rainbow
matching} problem. For each value $v$ define the menu
\[
E_v \;=\; \bigl\{\{v\}\bigr\}\;\cup\;\bigl\{\{x,v-x\} : 1\le x<v/2\bigr\},
\]
a set of vertex-disjoint edges (pairs summing to $v$ are pairwise disjoint) on the
vertex set $\intv{m_1}$; then $M$ is realizable iff one can select $\mu(v)$ pairwise
disjoint members of $E_v$ for every value $v$ simultaneously. (The singleton
members cannot be pre-committed here: the exchange behind
Corollary~\ref{cor:core} does not preserve the size cap, and $(6,6,5,4)$
witnesses the failure --- its stripped core demands a $6$ from $A=\{1,2,3\}$,
which no \emph{pair} provides, while in the pairs realization above the entry
$4$ takes $\{1,3\}$ rather than $\{4\}$. To reach the conventional edges-only
form, model the singleton $\{v\}$ instead as an edge from the vertex $v$ to a
private auxiliary vertex $s_v$ of color $v$; see the proof of
Proposition~\ref{prop:exactmatching}.) Full rainbow matching
is NP-hard in general, even under severe restrictions on the color classes: with all
classes matchings (properly edge-colored graphs) maximum rainbow matching is
APX-complete already on paths~\cite{LePfender14}; rainbow perfect matching is
NP-complete with color classes of size $\le2$~\cite{KelkStamoulis19}; and a recent
dichotomy~\cite{HommelsheimJM26} places ``every class a matching of size two'' on
the hard side. (The classical source of such hardness is the restricted perfect
matching problem of Itai, Rodeh and Tanimoto~\cite{ItaiRodehTanimoto78}; see the
attribution correction in~\cite{deFigueiredoMSS22}.) None of these instances have
the \emph{interval sum structure} of the menus $E_v$; whether that structure forces
tractability is, we believe, the right question:

\begin{question}\label{q:rainbow}
Is there a polynomial-time algorithm deciding, given multiplicities
$\mu:V\to\N$ on values $V\subseteq\intv{n}$, whether one can select $\mu(v)$
pairwise disjoint members of $E_v$ for each $v$ --- equivalently, deciding
realizability?
\end{question}

\subsubsection*{Status of the pairs conjecture}

Conjecture~\ref{conj:pairs} inherits the structural reductions of
Sections~\ref{sec:heredity} and~\ref{sec:minimal}, which localize any
possible counterexample.

\begin{proposition}[Pairs reductions]\label{prop:pairsred}
Call a realization with all $|X_i|\le2$ a \emph{pairs realization}, and $M$
\emph{pairs-realizable} if it admits one.
\begin{enumerate}[label=(\alph*),itemsep=1pt]
\item Pairs-realizability is hereditary under subsequences.
\item If $m_1>m_2$, then $M$ is pairs-realizable iff $T_2$ is; hence the
First-Repeat Reduction (Theorem~\ref{thm:firstrepeat}) holds verbatim under
the size cap.
\item If Conjecture~\ref{conj:pairs} fails, then some sequence is
simultaneously realizable and \emph{minimal} non-pairs-realizable (every
proper subsequence pairs-realizable); every such sequence is a
counterexample to the conjecture and has $m_1=m_2$.
\end{enumerate}
\end{proposition}

\begin{proof}
(a) As in Corollary~\ref{cor:heredity}: the kept blocks lie in the intrinsic
ground set by Lemma~\ref{lem:selfbound}, and the size cap is preserved.
(b) ($\Leftarrow$) A pairs realization of $T_2$ lies in $\intv{m_2}$
(Lemma~\ref{lem:selfbound}), and $X_1=\{m_1\}$ avoids it; ($\Rightarrow$) is
(a). (c) By (a) the non-pairs-realizable sequences form an up-set in the
subsequence order, so a realizable non-pairs-realizable $M$ contains a
minimal non-pairs-realizable $S$, and $S$ is realizable by
Corollary~\ref{cor:heredity} --- hence also a counterexample. If $S$ had
$m_1>m_2$, then by (b) its tail $T_2$ would be a proper
non-pairs-realizable subsequence, contradicting minimality.
\end{proof}

By Proposition~\ref{prop:pairsred}(c), verifying the conjecture up to a given
maximum requires only the realizable sequences with \emph{doubled maximum},
and by Theorem~\ref{thm:collapse} a DFS over realizable prefixes starting at
$(n,n)$ enumerates exactly these. This extends the verified range well
beyond $m_1\le12$ (\texttt{pairs\_scan.py}): for $13\le n\le18$ every
realizable sequence with $m_1=m_2=n$ --- there are $20{,}316$, $53{,}102$,
$151{,}214$, $405{,}982$, $1{,}171{,}102$ and $3{,}151{,}022$ of them,
$4{,}952{,}738$ in all --- was re-solved under the size cap, with no
exception. Hence \emph{Conjecture~\ref{conj:pairs} holds for every sequence
with $m_1\le18$}. Shorter sequences can be pushed much further: the
conjecture also holds for all sequences of length $k\le5$ with $m_1\le120$
($9{,}365{,}570$ realizable doubled-max instances), for all of length
$k\le6$ with $m_1\le40$ ($1{,}190{,}471$), and for all
$M=(n,n,a^{\alpha},b^{\beta})$ with $n\le26$ --- the two-repeated-value
shape of every forced-triple mechanism exhibited in this paper
(Theorem~\ref{thm:extremal}, the core of $(6,6,5,4)$, the Hall escapes of
Section~\ref{subsec:hall}). A minimal counterexample must therefore have
$m_1=m_2=n\ge19$, length $k\ge7$ if $n\le40$ and $k\ge6$ if $n\le120$, and,
for $19\le n\le26$, either a maximum of multiplicity at least three or at
least three distinct values below the maximum.

Beyond enumeration, the conjecture is a theorem in three regimes.

\begin{proposition}\label{prop:pairsspecial}
Every realizable $M$ is pairs-realizable if (i) the entries of $M$ are
pairwise distinct, (ii) $M$ is constant, or (iii) $M$ has length $k\le3$.
\end{proposition}

\begin{proof}
(i) All singletons (Corollary~\ref{cor:distinct}). (ii) By Ando, Gervacio
and Kano~\cite{AndoGervacioKano90}, $(m^{k})$ is realizable iff $m\ge2k-1$,
and exactly then $E_m$ contains $k$ pairwise disjoint members
$\{m\},\{1,m-1\},\dots,\{k-1,m-k+1\}$. (iii) By
Proposition~\ref{prop:pairsred}(b) assume $m_1=m_2=m$; the cases $(m,m)$ and
$(m,m,m)$ are covered by (ii). For $(m,m,b)$ with $b<m$ and $m\ge5$, take
$\{m\}$, $\{b\}$, and one of the disjoint pairs $\{1,m-1\},\{2,m-2\}$: the
single element $b$ meets at most one of them. For $m=4$ the only realizable
case is $(4,4,2)=\{4\},\{1,3\},\{2\}$, and for $m\le3$ no length-$3$
sequence with $m_1=m_2$ is realizable (Table~\ref{tab:catalog}).
\end{proof}

Two natural routes to a \emph{dis}proof close automatically. On the boundary
$\sum_im_i=\binom{n+1}{2}$ ($n=m_1$), where realizations of either kind
partition $\intv n$ (Section~\ref{sec:minimal}), a pairs realization
consists of exactly $2k-n$ singletons --- whose elements are their values,
hence pairwise distinct --- and $n-k$ pairs. Both resulting necessary
conditions already follow from realizability alone: any partition of
$\intv n$ into $k$ blocks has at least $2k-n$ singleton blocks with
pairwise distinct values, so the number of distinct values of a realizable
boundary sequence is at least $2k-n$; and since a block of odd (even) value
contains an odd (even) number of odd elements, any realization forces
$\#\{i:m_i\text{ odd}\}\equiv\lceil n/2\rceil\pmod2$ --- precisely the
parity a singleton--pair partition needs. More generally no element
exceeding $n/2$ can separate the two notions: no block of either kind
contains two such elements, and a block containing $e$ has value at least
$e$. A counterexample must be blocked through the \emph{small} elements, as
a genuine matching failure invisible to counting.

On the proof side, the obstacle reduces to orienting one exchange. Among
realizations of $M$, minimize $\Phi=\sum_{|X_i|\ge3}|X_i|$ and then
$\Phi'=\sum_{|X_i|\ge3}m_i$ lexicographically; $\Phi=0$ iff the realization is
a pairs realization, so in a counterexample to Conjecture~\ref{conj:pairs}
every $(\Phi,\Phi')$-minimal realization has a block of size $\ge3$. Such
blocks are severely constrained:

\begin{proposition}[Pair saturation]\label{prop:saturation}
In a $(\Phi,\Phi')$-minimal realization of $M$, every block $X$ of size
$\ge3$ has value $v<m_1$ and is \emph{upward pair-saturated}: for each
$a\in X$ with $2a\neq v$, the element $v-a$ lies in a pair of value strictly
greater than $v$ whose other element exceeds $a$.
\end{proposition}

\begin{proof}
Let $X$ be a block of size $\ge3$ and value $v$, let $a\in X$ with $2a\neq v$
(such an $a$ exists: at most one element of $X$ equals $v/2$), and set
$w=v-a$. Then $w\notin X$ --- otherwise $a$, $w$ and a third positive element
would give $\mass(X)>v$ --- and every element of $X\setminus\{a\}$ is smaller
than $w=\mass(X\setminus\{a\})$, that set having at least two elements. We
locate the element $w$ in the realization and exchange.

If $w$ is unused, replace $X$ by $\{a,w\}$: $\Phi$ strictly decreases. If $w$
lies in a singleton block --- whose entry then has value $w$ --- give that
entry the set $X\setminus\{a\}$, of the same mass, and give $X$'s entry
$\{a,w\}$: again $\Phi$ strictly decreases. If $w$ lies in a block $Y$ of
size $\ge3$, replace $Y$ by $(Y\setminus\{w\})\cup(X\setminus\{a\})$ (the
mass is unchanged) and $X$ by $\{a,w\}$: $\Phi$ decreases by $2$. Each move
preserves disjointness, since $X\cap Y=\emptyset$ and $w\notin X$; each
therefore contradicts minimality.

Hence $w$ lies in a pair $\{w,u\}$, and the swap
$(X,\{w,u\})\mapsto(\{a,w\},\,\{u\}\cup X\setminus\{a\})$ is available: both
entries keep their values ($a+w=v$, and $u+\mass(X\setminus\{a\})=u+w$),
$\Phi$ is unchanged, and $\Phi'$ changes by $u-a$. Minimality forces $u>a$,
so the pair $\{w,u\}$ has value $w+u>w+a=v$: $X$ is upward pair-saturated.
Finally, if $v=m_1$ then $w+u\le m_1=w+a$ gives $u\le a$, a contradiction; so
no block of size $\ge3$ has value $m_1$.
\end{proof}

Any potential function monotone under the pair swap would prove
Conjecture~\ref{conj:pairs}; together with the deficiency form of
Question~\ref{q:rainbow}, Proposition~\ref{prop:saturation} seems to us the
sharpest formulation of the conjecture's difficulty.

The exact layered solver behind every computational claim in this paper --- complete
independently of Conjecture~\ref{conj:pairs} --- is described in
Appendix~\ref{app:solver}.

\section{Minimal forbidden sequences}\label{sec:minimal}

Forbidden sequences form an up-set in the subsequence order
(Section~\ref{sec:heredity}), so they are generated by their minimal elements.
This section studies the generators: how repetition constrains their shape, the
complete catalog for $m_1\le12$, an extremal family answering how large their
entries can be, and a boundary theorem locating them strictly inside the volume
region.

\subsection{Minimality and the first-repeat reduction}

\begin{proposition}\label{prop:minimalwitness}
$M$ is forbidden iff $M$ contains a minimal forbidden subsequence (forbidden, with
every proper subsequence realizable). Testing minimality requires only the $k$
single-deletion subsequences.
\end{proposition}

\begin{proof}
($\Leftarrow$) Corollary~\ref{cor:heredity}. ($\Rightarrow$) Among forbidden
subsequences of $M$ pick one of minimum length. For the second claim: every proper
subsequence lies below some single-deletion subsequence in the subsequence order,
and realizability is downward closed.
\end{proof}

\begin{observation}\label{obs:doubledmax}
If $m_1>m_2$ then $M$ is realizable iff $T_2$ is: take $X_1=\{m_1\}$, disjoint from
all other blocks by Lemma~\ref{lem:selfbound}. Hence every minimal forbidden
sequence has $m_1=m_2$.
\end{observation}

Iterating gives a stronger reduction --- forbiddenness can jump straight to the first
repeated entry:

\begin{theorem}[First-Repeat Reduction]\label{thm:firstrepeat}
Let $i^*$ be the least index with $m_{i^*}=m_{i^*+1}$ (if no entry repeats, $M$ is
realizable by Corollary~\ref{cor:distinct}). Then $M$ is realizable iff $T_{i^*}$ is.
Equivalently, realizability depends only on the entries at most the largest repeated
value.
\end{theorem}

\begin{proof}
While the current maximum is strictly larger than the next entry, drop it
(Observation~\ref{obs:doubledmax}); the process stops exactly at $T_{i^*}$.
\end{proof}

The reduction \emph{point} is the first repeat, but the minimal \emph{witness} may
sit deeper: $(9,9,2,2)$ has first repeat at $9$, yet $(9,9)$ is realizable and the
minimal witness is $(2,2)$.

\subsection{The catalog for $m_1\le12$}

Minimal forbidden sequences can be enumerated with no length bound: extensions of
forbidden sequences are never minimal, so a DFS over \emph{realizable} prefixes
starting $(n,n)$ visits exactly the candidates. (The minimal forbidden sequences of
length at most $3$ --- the rows $n\le4$ of Table~\ref{tab:catalog} --- were already
determined in~\cite[\S4.2.1]{Moragas10}.)

\begin{table}[ht]
\centering
\caption{Minimal forbidden sequences by maximum entry $n\le12$
(\texttt{minimal\_forbidden.py}). The count sequence is not in the OEIS as of this
writing.}
\label{tab:catalog}
\begin{tabular}{@{}rrll@{}}
\toprule
$n$ & count & lengths & examples (shortest) \\
\midrule
1 & 1 & 2 & $(1,1)$ \\
2 & 1 & 2 & $(2,2)$ \\
3 & 3 & 3 & $(3,3,1)$, $(3,3,2)$, $(3,3,3)$ \\
4 & 3 & 3 & $(4,4,1)$, $(4,4,3)$, $(4,4,4)$ \\
5 & 11 & 4 & $(5,5,2,1)$, $(5,5,3,1)$, \dots \\
6 & 12 & 4--5 & $(6,6,2,1)$, \dots, $(6,6,5,4,3)$ \\
7 & 42 & 4--5 & $(7,7,3,3)$ unique of length 4 \\
8 & 57 & 4--6 & $(8,8,3,3)$ unique of length 4 \\
9 & 158 & 5--7 & $(9,9,9,3,3)$, $(9,9,6,6,2)$ \\
10 & 291 & 5--8 & $(10,10,10,3,3)$, $(10,10,4,4,2)$ \\
11 & 734 & 5--8 & $(11,11,5,5,5)$ unique of length 5 \\
12 & 1355 & 5--9 & $(12,12,5,5,5)$ unique of length 5 \\
\bottomrule
\end{tabular}
\end{table}

Structural statistics over the $2{,}668$ sequences of Table~\ref{tab:catalog}: the
maximum entry has multiplicity $2$ in $1{,}406$ of them but repeats up to $7$ times;
$1{,}995$ have at least two distinct repeated values (so
Observation~\ref{obs:doubledmax} does not exhaust the repetition structure); and
$2{,}443$ contain at least one unrepeated entry --- unrepeated values matter because
stripping deletes their value from the ground set, as in $(6,6,5,4,3)$, where the
distinct tail $5,4,3$ leaves $A=\{1,2\}$.

\subsection{An extremal family and the length frontier}

How large can the maximum entry of a length-$k$ minimal forbidden sequence be? Let
$n_{\max}(k)$ denote the maximum. The question was raised in Moragas's
thesis~\cite[\S4.2.1]{Moragas10}, where the length-$k$ sequence
$(2a{+}1,\,2a{+}1,\,a,\dots,a)$, with $k-2$ copies of $a=2k-5$, is shown to be
minimal forbidden --- giving $n_{\max}(k)\ge4k-9$ --- and the upper bound
$n_{\max}(k)<4k$ is conjectured. The following family improves the lower bound by
one; we conjecture that its value is exact.

\begin{theorem}[Extremal family]\label{thm:extremal}
For every $j\ge1$ the sequence
$F_j=(4j,\,4j,\,\underbrace{2j-1,\dots,2j-1}_{j\text{ copies}})$, of length $j+2$,
is minimal forbidden. Hence $n_{\max}(k)\ge4(k-2)$ for $k\ge3$.
\end{theorem}

\begin{proof}
\emph{Forbidden.} Strip the singletons $\{4j\}$ and $\{2j-1\}$
(Corollary~\ref{cor:core}). The $j-1$ excess copies of $2j-1$ demand disjoint
subsets of $\intv{2j-2}$ (each element is $\le 2j-2$: the value $2j-1$ is stripped)
whose total, $(j-1)(2j-1)$, equals $\mass(\intv{2j-2})$ exactly --- so any
realization must partition $\intv{2j-2}$, which is possible via the pairs
$\{x,\,2j-1-x\}$. Everything at most $2j-1$ is then consumed, and the excess copy of
$4j$ must be drawn from $\{2j,\dots,4j-1\}$: a single element is at most $4j-1$, and
the two cheapest available elements sum to $(2j)+(2j+1)=4j+1>4j$. Blocked.

\emph{Minimal.} Deleting one $4j$ leaves $(4j,(2j-1)^{\times j})$: realize by
$\{4j\}$, $\{2j-1\}$, and the pairs partition of $\intv{2j-2}$. For the deletion of
one $2j-1$ the case $j=1$ is degenerate --- the pairs partition of
$\intv{2j-2}=\emptyset$ contains no pair to free --- and is checked directly:
deleting the $1$ from $F_1=(4,4,1)$ leaves $(4,4)$, realized by $\{4\},\{1,3\}$.
For $j\ge2$, deleting one $2j-1$ frees one pair $\{y,\,2j-1-y\}$ of the partition,
and the second $4j$ is realized as $\{y,\,2j-1-y,\,2j+1\}$, of sum
$(2j-1)+(2j+1)=4j$. All single-deletions are realizable, hence
(Proposition~\ref{prop:minimalwitness}) all proper subsequences are.
\end{proof}

The family begins $(4,4,1)$, $(8,8,3,3)$, $(12,12,5,5,5)$, $(16,16,7,7,7,7)$ --- its
first two members are exactly the volume-passing examples after
Theorem~\ref{thm:collapse}.

\begin{conjecture}[Frontier]\label{conj:frontier}
$n_{\max}(k)=4(k-2)$ for all $k\ge3$; for $k\ge4$ the extremal witness is unique,
namely $F_{k-2}$.
\end{conjecture}

Verified by exhaustive scan: $k=3$ up to $n=12$ ($n_{\max}=4$, witnesses
$(4,4,1),(4,4,3),(4,4,4)$); $k=4$ up to $n=20$ (unique witness $(8,8,3,3)$); $k=5$
up to $n=24$ (unique $(12,12,5,5,5)$); $k=6$ up to $n=18$ (unique
$(16,16,7,7,7,7)$). The contrapositive reading explains the shape: blocking a large
value $n$ requires consuming an entire initial segment $\intv t$ with
$t\ge(n-2)/2$, and forcing that consumption takes on the order of $t/2$ coupled
entries.

\subsection{No minimal forbidden sequence on the partition boundary}

The exact-partition regime of the sumset partition problem corresponds, inside the
relaxed problem, to the boundary case
$\sum_i m_i=\binom{m_1+1}{2}$ of the volume inequality: there, disjointness plus
full weight force any realization to partition $\intv{m_1}$, so
\emph{realizable $\Leftrightarrow$ $m_1$-realizable} on the boundary.

\begin{lemma}[Boundary complement]\label{lem:boundarycomp}
Let $M$ satisfy $\sum_im_i=\binom{m_1+1}{2}$ with $k\ge2$, and let $M'$ be $M$ with
one entry $m_j$ deleted. If $M'$ is realizable, so is $M$.
\end{lemma}

\begin{proof}
Take a realization $\{X_i\}_{i\neq j}$ of $M'$ inside
$\intv{\max M'}\subseteq\intv{m_1}$ and let
$U=\intv{m_1}\setminus\bigcup_{i\neq j}X_i$. Then
\[
\mass(U)\;=\;\binom{m_1+1}{2}-\sum_{i\neq j}m_i\;=\;m_j
\]
by the boundary equality, so $X_j:=U$ completes a realization of $M$.
\end{proof}

\begin{theorem}\label{thm:boundary}
No minimal forbidden sequence satisfies $\sum_im_i=\binom{m_1+1}{2}$. Equivalently:
on the exact-partition boundary, forbiddenness --- when present --- always comes from
a proper subsequence (necessarily of positive volume slack).
\end{theorem}

\begin{proof}
For $k=1$ the boundary forces $m_1=1$ and $(1)$ is realizable. For $k\ge2$, if $M$
is forbidden on the boundary then by Lemma~\ref{lem:boundarycomp} (contrapositive)
\emph{every} single-deletion subsequence is forbidden, so $M$ is not minimal.
\end{proof}

Three remarks. (i) Combining Lemma~\ref{lem:boundarycomp} with heredity, a boundary
sequence with $k\ge2$ is realizable iff any one (hence every) single-deletion
subsequence is realizable; choosing a deletion with $j\ge2$ (possible as $k\ge2$)
keeps the maximum, so $M'$ has volume slack $m_j>0$: the equivalence takes one step
off the boundary and stops --- which is exactly why minimal witnesses live at
positive slack. (Deleting the maximum can instead shrink the ground set and land
back on a boundary, as with $M=(2,1)\mapsto(1)$.) (ii) The Lemma's construction converts any realization of $M'$ into
an explicit partition of $\intv{m_1}$. (iii) The $n$-feasible-but-not-$n$-realizable
sequences of the sumset partition problem are precisely the boundary sequences
containing an off-boundary minimal witness (Proposition~\ref{prop:minimalwitness}
plus Theorem~\ref{thm:boundary}).

\section{The assignment tree: $n$-realizability through the relaxed
oracle}\label{sec:tree}

We now return to the exact problem. Fix an $n$-feasible $M$. The \emph{assignment
tree} $T(M,n)$ has root $(M,n)$; a node $(S,m)$ holds an $m$-feasible non-increasing
sequence $S$ with \emph{pending element} $m$. Its children are obtained by assigning
the element $m$ to an entry of value $v\ge m$: replace one copy of $v$ by $v-m$
(deleting it if $v=m$) and re-sort. Children are $(m-1)$-feasible. Assigning $m$ to
equal entries gives the same child, so a node has as many children as distinct
values $\ge m$. The empty sequence appears exactly at pending element $0$, so a
maximal root-to-leaf path either reaches the \emph{empty leaf} at depth $n$ or stops
early at a \emph{dead end}: a nonempty node with no entry $\ge m$.

\begin{lemma}\label{lem:deadend}
Dead ends are forbidden (in the relaxed sense).
\end{lemma}

\begin{proof}
At a dead end $(S,m)$ every entry is $\le m-1$, so $s_1=\max S\le m-1$, while
$\sum S=\binom{m+1}{2}>\binom{m}{2}\ge\binom{s_1+1}{2}$: the $i=1$ volume
condition~\eqref{eq:volume} fails.
\end{proof}

\begin{theorem}[Assignment-tree characterization]\label{thm:tree}
For $n$-feasible $M$ the following are equivalent:
\begin{enumerate}[label=(\alph*),itemsep=1pt]
\item $M$ is $n$-realizable;
\item some maximal path of $T(M,n)$ has all its nodes realizable (relaxed sense);
\item some path of $T(M,n)$ reaches the empty leaf.
\end{enumerate}
\end{theorem}

\begin{proof}
(c)$\Rightarrow$(a): a root-to-empty path assigns each element $n,n-1,\dots,1$ to a
block, and reaching the empty leaf means every block's residual hit zero: the
assignments are a partition of $\intv n$ with block sums $M$.
(a)$\Rightarrow$(b): fix a partition $X_1,\dots,X_k$ and follow the path assigning
each element $m$ to the block containing it (that block's residual is a sum of
distinct unassigned elements including $m$, hence $\ge m$). The node at pending $m$
holds the residuals $r_i=\mass(X_i\cap\intv m)$, realized by the sets
$X_i\cap\intv m$; by Lemma~\ref{lem:selfbound} these lie in the node's own ground
set, so the node is realizable.
(b)$\Rightarrow$(c): by Lemma~\ref{lem:deadend} a maximal path with no forbidden
node cannot end at a dead end.
\end{proof}

Thus tree search with relaxed-oracle pruning is a \emph{correct and complete}
decision procedure for $n$-realizability: an all-realizable maximal path \emph{is} a
partition, and pruning at forbidden nodes is sound. Unlike
Theorem~\ref{thm:collapse}, where the analogous reduction collapsed to a single
test, this recursion has genuine depth --- quantified next.

\section{The order hierarchy}\label{sec:order}

Theorem~\ref{thm:tree} makes the relaxed oracle a sound and complete pruning
rule for the exact problem; the depth at which the pruning fires measures,
instance by instance, how much of $n$-realizability the relaxation fails to
see. This section quantifies that depth.

\begin{definition}\label{def:order}
For $n$-feasible, non-$n$-realizable $M$, every maximal path of $T(M,n)$ contains a
forbidden node (Theorem~\ref{thm:tree}). Define $\ord(M)$ as the maximum over
maximal paths of the level of the first forbidden node (root $=$ level 1).
Equivalently: $\ord(M)=1$ if $M$ is forbidden; otherwise
$\ord(M)=1+\max_{\text{children }C}\ord(C)$. Here all children of a
non-$n$-realizable node are non-realizable --- if the child obtained by assigning
the pending element $m$ to an entry of value $v$ had a partition of $\intv{m-1}$,
then adding the element $m$ to the block realizing $v-m$ (or adding the new block
$\{m\}$ when $v=m$) would $m$-realize the node --- and a realizable node has at
least one child, since \eqref{eq:volume} gives $m_1\ge n$. Let $\maxord(n)$ be the maximum order over all
$n$-feasible non-$n$-realizable sequences.
\end{definition}

Order $1$ means the relaxed oracle refutes $M$ at the root. Order $\ge2$ means $M$
is realizable in the relaxed sense yet not $n$-realizable: the obstruction is
invisible to the relaxation until at least one assignment has been made. The order
is precisely the refutation depth of the pruned tree search --- a per-instance
hardness measure for the relaxation.

\subsection{Lower-order structure: swap and second-entry bounds}

\begin{theorem}[One-overflow swap]\label{thm:swap}
If an $n$-feasible $M$ has a relaxed realization using at most one element $>n$,
then $M$ is $n$-realizable. Hence every sequence of order $\ge2$ has $m_1\ge n+2$,
and every one of its realizations uses at least two elements exceeding $n$.
\end{theorem}

\begin{proof}
Let $\{X_i\}$ be such a realization and $U=\intv n\setminus\bigcup_iX_i$. The used
mass is $\binom{n+1}{2}$, so $\mass(U)$ equals the used mass above $n$. If no
element $>n$ is used, the $X_i$ partition $\intv n$. If exactly one element $e>n$ is
used, say $e\in X_j$, then $\mass(U)=e$: replace $X_j$ by
$(X_j\setminus\{e\})\cup U$ --- same sum, still disjoint --- and the previous case
applies. With $m_1\le n+1$ at most one overflow element exists at all.
\end{proof}

\begin{theorem}[Second-entry bound]\label{thm:secondentry}
If $M$ is $n$-feasible, realizable, and $m_2\le n$, then $M$ is $n$-realizable.
Consequently every order-$\ge2$ sequence has $m_1\ge n+2$ \emph{and} $m_2\ge n+1$.
\end{theorem}

\begin{proof}
By heredity the tail $(m_2,\dots,m_k)$ is realizable with blocks
$X_i\subseteq\intv{m_i}\subseteq\intv n$ (Lemma~\ref{lem:selfbound}). Let
$U=\intv n\setminus\bigcup_{i\ge2}X_i$; then
$\mass(U)=\binom{n+1}{2}-(\sum_im_i-m_1)=m_1$ by feasibility, so $X_1:=U$ completes
a partition. Note no assumption is made on how $m_1$ was realized: the first block
is rebuilt from the leftovers.
\end{proof}

Theorem~\ref{thm:secondentry} constrains the tree: an order-$\ge2$ node always has
at least two entries exceeding its pending element, so either its top two entries
are equal (one child) or it has $\ge2$ children --- deep refutations never ride a
single forced line unless the maximum is doubled.

Order-$\ge2$ sequences exist from $n=6$ on, and $m_1\ge n+2$ is sharp. The three
smallest are $(8,8,3,2)$, $(8,7,3,3)$, $(8,7,3,2,1)$ at $n=6$. For $M=(8,8,3,2)$: in
$\intv6$ the blocks $2=\{2\}$ and $3=\{3\}$ are forced ($\{1,2\}$ would strand
$1$), leaving $\{1,4,5,6\}$ with no subset summing to $8$ --- not $6$-realizable;
but $\{8\},\{1,7\},\{3\},\{2\}$ is a relaxed realization using the two overflow
elements $7,8$. Its only child, $(8,3,2,2)$, is forbidden (the excess $2$ has no
partner), so $\ord(8,8,3,2)=2$. The third witness, $(8,7,3,2,1)$, has pairwise
\emph{distinct} entries; such sequences are the subject of the next subsection.

\subsection{Distinct entries: counting obstructions and counting-blind
witnesses}\label{subsec:distinct}

By Corollary~\ref{cor:distinct} a sequence with pairwise distinct entries is always
realizable in the relaxed sense, so every distinct-entry $n$-feasible
non-$n$-realizable sequence has order $\ge2$ automatically: with distinct entries,
non-realizability is a pure \emph{coverage} obstruction --- exactly the large-number
placement that the relaxation discards. (An early example is $(13,12,11,9,4,3,2,1)$
at $n=10$, noted in~\cite[Example~4.1]{Moragas10} precisely to show that the absence
of forbidden subsequences does not imply $n$-realizability.) The witness
$(8,7,3,2,1)$ suggests where such obstructions become visible: not in the
relaxation, but in counting elements. Each block of a partition of $\intv n$ is a
set of distinct integers from $\intv n$, which bounds the number of its elements
exceeding any threshold from both sides in terms of its sum.

\begin{proposition}[Counting conditions]\label{prop:counting}
For integers $0\le t<n$ and $m\ge1$ put
\[
\alpha_t(m)=\max\Bigl\{j\ge0:\ jt+\tbinom{j+1}{2}\le m\Bigr\},
\qquad
\lambda_t(m)=\min\Bigl\{j\ge0:\ jn-\tbinom{j}{2}\,\ge\, m-\tbinom{t+1}{2}\Bigr\}
\]
(the latter well defined whenever $m\le\binom{n+1}{2}$, as is every entry of an
$n$-feasible sequence). If $M$ is $n$-realizable, then
\[
\sum_{i=1}^{k}\lambda_t(m_i)\ \le\ n-t\ \le\ \sum_{i=1}^{k}\alpha_t(m_i)
\qquad\text{for every } t\in\{0,1,\dots,n-1\}.
\]
\end{proposition}

\begin{proof}
Fix a partition $X_1\cup\dots\cup X_k=\intv n$ realizing $M$ and let
$j_i=|X_i\cap\{t+1,\dots,n\}|$, so that $\sum_ij_i=n-t$ exactly. The $j_i$ elements
of $X_i$ exceeding $t$ are distinct integers $\ge t+1$, so they carry mass at least
$j_it+\binom{j_i+1}{2}$, which is at most $m_i$; hence $j_i\le\alpha_t(m_i)$. The
elements of $X_i$ not exceeding $t$ are distinct positive integers $\le t$, of total
mass at most $\binom{t+1}{2}$, so the elements exceeding $t$ carry at least
$m_i-\binom{t+1}{2}$; being $j_i$ distinct integers $\le n$, they carry at most
$j_in-\binom{j_i}{2}$, whence $j_i\ge\lambda_t(m_i)$. Summing over $i$ gives both
inequalities.
\end{proof}

At $t=0$ the proposition reads $\sum_i\lambda_0(m_i)\le n\le\sum_i\alpha_0(m_i)$:
the blocks jointly need at least $\sum_i\lambda_0(m_i)$ elements and can hold at
most $\sum_i\alpha_0(m_i)$, while a partition uses exactly $n$. Unlike the
Hall-type conditions of Section~\ref{subsec:hall}, which constrain the relaxed core
inside $\intv{m_1}$, these conditions have no relaxed analogue: both inequalities
rest on the covering constraint $\sum_ij_i=n-t$, precisely the information the
relaxation discards. For $(8,7,3,2,1)$ at $n=6$ the count fails outright: the
entries $8$ and $7$ exceed $6$, so each needs at least two elements of $\intv6$,
while $3$, $2$ and $1$ need one each --- at least seven elements, and only six
exist. This is the first member of an infinite family.

\begin{proposition}[A distinct order-$2$ family at every level]\label{prop:Dn}
For $n\ge6$ let
\[
D_n=(n+2,\ n+1,\ n-2,\ n-3,\ \dots,\ 5,\ 3,\ 2,\ 1),
\]
the interval $\intv{n-2}$ with the entry $4$ removed, preceded by the two towers
$n+2$ and $n+1$. Then $D_n$ is $n$-feasible with pairwise distinct entries, violates
the count condition ($t=0$) of Proposition~\ref{prop:counting}, and
$\ord(D_n)=2$. In particular the bounds $m_1\ge n+2$ and $m_2\ge n+1$ of
Theorems~\ref{thm:swap} and~\ref{thm:secondentry} are attained simultaneously at
every level $n\ge6$; $D_6=(8,7,3,2,1)$.
\end{proposition}

\begin{proof}
Feasibility: $\sum D_n=\binom{n-1}{2}-4+(2n+3)=\binom{n-1}{2}+(2n-1)
=\binom{n+1}{2}$; the entries are distinct and sorted for $n\ge6$. Counting: the
$n-3$ entries below $n$ have $\lambda_0=1$, while the towers exceed $n$ and are at
most $2n-1$, so $\lambda_0(n+2)=\lambda_0(n+1)=2$; the total is $n+1>n$, so $D_n$ is
not $n$-realizable. Distinct entries make $D_n$ realizable
(Corollary~\ref{cor:distinct}), so $\ord(D_n)\ge2$. For the exact value: the pending
element $n$ exceeds every entry except the two towers, so the root has exactly two
children, obtained by assigning $n$ to $n+2$ (residual $2$) or to $n+1$ (residual
$1$). Since the tail already contains the entries $2$ and $1$, the children contain
the minimal forbidden sequences $(2,2)$ and $(1,1)$ respectively
(Table~\ref{tab:catalog}) and are forbidden by Theorem~\ref{thm:collapse}. Hence
every maximal path meets a forbidden node at level $2$, and $\ord(D_n)=2$.
\end{proof}

\begin{remark}[The two oracles are incomparable]\label{rem:incomparable}
Neither refutation engine subsumes the other. The floor-and-towers sequence $W_a$ of
Theorem~\ref{thm:Wa} below has order $a$ --- the relaxed oracle needs depth $a$ to
refute it --- yet the count condition kills it at the root: at $n=2a+1$ the $a-1$
largest elements have mass $\tfrac{3a^2+a-4}{2}<P$, so $\lambda_0(P)\ge a$, while
$\lambda_0(Q)=2$ and the floor contributes $a$, for a total of at least
$2a+2=n+1>n$; the
same estimate applies to the even-level family $V_r$, and
Moragas's example $(13,12,11,9,4,3,2,1)$ fares no better ($\lambda_0$ sums to
$11>10$). Conversely $\bigl(\binom{n+1}{2}-6,\,4,\,1,\,1\bigr)$ contains the minimal
forbidden $(1,1)$ --- order $1$ --- yet satisfies every condition of
Proposition~\ref{prop:counting} (verified for $n\le400$). Pruning the assignment
tree at nodes failing either oracle is of course still sound, and can only decrease
the refutation depth.
\end{remark}

So counting refutes $D_n$, $W_a$ and $V_r$ at the root, and the natural question is
whether a coverage obstruction can hide from counting entirely. It can: the smallest
counting-clean distinct-entry witnesses are $(12,8,5,2,1)$ and $(10,8,6,3,1)$ at
$n=7$ (census counts in Section~\ref{subsec:census}), and an infinite family rides a
parity mechanism instead.

\begin{theorem}[Odd towers]\label{thm:oddtowers}
For even $n=2s\ge8$ let
\[
O_n=\bigl(u,\ v,\ n-1,\ n-3,\ \dots,\ 3,\ 1\bigr),\qquad
u=s^2-s-1=\tfrac{n^2-2n-4}{4},\quad v=n+1,
\]
the two towers $u>v$ followed by every odd number below $n$ --- all entries odd.
Then:
\begin{enumerate}[label=(\alph*),itemsep=1pt]
\item $O_n$ is $n$-feasible with pairwise distinct entries;
\item $O_n$ satisfies every counting condition of Proposition~\ref{prop:counting};
\item $O_n$ is not $n$-realizable;
\item $\lfloor n/4\rfloor\le\ord(O_n)\le n-5$.
\end{enumerate}
\end{theorem}

\begin{proof}
(a) The floor has mass $\sum_{k=1}^{s}(2k-1)=s^2$ and $u+v=s^2+s$, so
$\sum O_n=2s^2+s=\binom{n+1}{2}$. All entries are odd, and $u>v$ amounts to
$s^2-3s-2>0$, true for $s\ge4$.

(c) In any partition of $\intv n$ realizing $O_n$, the block of the entry $1$ is
$\{1\}$. Inductively, suppose the entries $1,3,\dots,2k-3$ hold the blocks
$\{1\},\{3\},\dots,\{2k-3\}$. The block of the entry $2k-1$ has odd sum, hence
contains an odd element; its elements are at most $2k-1$ (Lemma~\ref{lem:selfbound}),
and every smaller odd element is committed; so it contains the element $2k-1$ and,
by mass, equals $\{2k-1\}$. The floor therefore consumes exactly the odd elements of
$\intv n$, and the blocks of $u$ and $v$ would have to partition the even elements
$\{2,4,\dots,n\}$ --- impossible, since a set of even numbers has even mass while $u$
and $v$ are odd. (No global parity count is violated: as for any $n$-feasible
sequence, $\#\{i:m_i\text{ odd}\}=s+2\equiv\lceil n/2\rceil\pmod 2$; the obstruction
is local, appearing only after the floor is forced.)

(b) Write $\mathrm{top}_j=jn-\binom{j}{2}$ for the mass of the $j$ largest elements
of $\intv n$ and $\mathrm{bot}_j(t)=jt+\binom{j+1}{2}$ for that of the $j$ smallest
elements exceeding $t$, so that $\lambda_t(m)=\min\{j:\mathrm{top}_j\ge
m-\binom{t+1}{2}\}$ and $\alpha_t(m)=\max\{j:\mathrm{bot}_j(t)\le m\}$.

\emph{Lower conditions.} First, $\lambda_t(m)\le n-t$ for every single entry, since
$\mathrm{top}_{n-t}=\binom{n+1}{2}-\binom{t+1}{2}\ge m-\binom{t+1}{2}$. Second,
$\mathrm{top}_{s-2}=\tfrac32(s-2)(s+1)\ge s^2-s-1$ (equivalent to $s^2-s\ge4$), so
$\lambda_t(u)\le\lambda_0(u)\le s-2$ for every $t$. At $t=0$: $\lambda_0(v)=2$ and
each of the $s$ floor entries has $\lambda_0=1$, so the total is at most
$(s-2)+2+s=n$. For $1\le t\le s+1$: $\lambda_t(v)\le1$ (as
$v-\binom{t+1}{2}\le n=\mathrm{top}_1$), and a floor entry contributes only if
$2k-1>\binom{t+1}{2}$; since $\binom{t+1}{2}\ge2t-3$, at most $s-(t-1)$ values of
$k$ qualify, so the total is at most $(s-2)+1+(s-t+1)=n-t$. For $t\ge s+2$:
$\binom{t+1}{2}\ge\binom{s+3}{2}\ge2s+1$, so $v$ and the floor contribute nothing,
and the total is $\lambda_t(u)\le n-t$ by the per-entry bound.

\emph{Upper conditions.} A floor entry can absorb an element exceeding $t$ iff
$2k-1\ge t+1$, which happens for $s-\lceil t/2\rceil$ values of $k$; and
$\alpha_t(v)\ge2$ iff $2t+3\le2s+1$, i.e.\ iff $t\le s-1$ (always
$\alpha_t(v)\ge1$). It therefore suffices to show $\alpha_t(u)\ge J$, where
$J=s-\lfloor t/2\rfloor-2$ for $0\le t\le s-1$ and $J=s-\lfloor t/2\rfloor-1$ for
$s\le t\le 2s-1$ (nothing to prove when $J\le0$). Put $w=\lfloor t/2\rfloor$, so
$t\le2w+1$ and $\mathrm{bot}_J(t)\le J(2w+1)+\binom{J+1}{2}$. In the first case
$w=s-J-2$ and the bound equals $J(4s-3J-5)/2\le(4s-5)^2/24\le s^2-s-1=u$, the last
inequality being $8s^2+16s\ge49$. In the second case $w=s-J-1$ with
$J\le s-\lfloor s/2\rfloor-1\le(s-1)/2$, and the bound equals $J(4s-3J-1)/2$,
increasing in $J$ up to $(4s-1)/6\ge(s-1)/2$, hence at most
$(s-1)(5s+1)/8\le s^2-s-1=u$, the last inequality being $(3s-7)(s+1)\ge0$. In both
cases $\mathrm{bot}_J(t)\le u$, i.e.\ $\alpha_t(u)\ge J$, and the condition holds.

(d) For the lower bound, follow the path that assigns every pending element to the
$u$-tower. After $j$ assignments the tower's value is
$u_j=u-\sum_{i=n-j+1}^{n}i\ge u-jn$, so for $j\le(s-4)/2$ we get
$u_j\ge s^2-s-1-s(s-4)=3s-1>2s+1=v$: the node at level $j+1$ has entries
$u_j,v,n-1,\dots,3,1$, pairwise distinct, hence realizable
(Corollary~\ref{cor:distinct}), and the next assignment is legal
($u_j>n\ge{}$pending). Some path is thus alive through level
$\lfloor(s-4)/2\rfloor+1$, so the first forbidden node on any of its maximal
extensions sits at level at least
$\lfloor(s-4)/2\rfloor+2=\lfloor s/2\rfloor=\lfloor n/4\rfloor$; the order,
being the maximum of that level over all maximal paths, is at least
$\lfloor n/4\rfloor$. The upper bound is Theorem~\ref{thm:upper}.
\end{proof}

Machine computation (\texttt{distinct\_families.py}) gives
$\ord(O_n)=2,2,4,6,8,8,10,12,16,18,20,22,24$ for $n=8,10,\dots,32$ --- equal to
$n-8$ for every even $24\le n\le32$, three below the maximum $\maxord(n)=n-5$
(Table~\ref{tab:lift}). Counting-blind distinct-entry witnesses of near-maximal
order thus exist as far as we can compute; whether $\ord(O_n)=n-8$ persists, and
whether distinct-entry (or counting-blind) sequences can attain $\maxord(n)$ itself,
we leave open (Section~\ref{sec:open}).

\subsection{Census for $n\le12$}\label{subsec:census}

All $n$-feasible sequences with $n\le12$ were classified and the order of every
non-$n$-realizable, realizable one computed (\texttt{assignment\_tree.py}); the full
counts, by order, appear in Table~\ref{tab:census} of Appendix~\ref{app:census}.
The fraction of realizable feasible sequences that are $n$-realizable declines
slowly with $n$, from $97\%$ at $n=6$ to $87\%$ at $n=12$, and the maximum order
climbs by exactly one with each $n$ from $n=8$ on. High order is built
by \emph{lifting}: the unique order-$5$ sequence at $n=10$ is
$(16,12,10,10,4,2,1)$; adding the element $11$ back onto the entry $1$ gives
$(16,12,12,10,10,4,2)$, one of three order-$6$ sequences at $n=11$; and the unique
order-$7$ sequence at $n=12$, $(24,16,11,10,10,4,2,1)$, has as one child the
order-$6$ witness $(16,12,11,10,10,4,2,1)$. This un-assignment mechanism is
formalized in Proposition~\ref{prop:lift}. Distinct-entry non-realizable sequences
are not rare: $1,6,20,73,215,692,1{,}972$ of them for $n=6,\dots,12$, of which
$0,2,6,31,102,389,1{,}166$ satisfy every counting condition of
Proposition~\ref{prop:counting}. The largest order attained by a distinct-entry
sequence is $2,3,3,4,4,5,6$ for $n=6,\dots,12$ --- equal to $\maxord(n)$ up to
$n=9$ (attained at $n=7$ by $W_3=(14,8,3,2,1)$ of Theorem~\ref{thm:Wa}), and
exactly one short of it for $10\le n\le12$ (\texttt{distinct\_families.py}).

\subsection{The upper bound $\maxord(n)\le n-5$}

\begin{theorem}[Step bound]\label{thm:step}
$\maxord(n+1)\le1+\maxord(n)$ for all $n\ge1$.
\end{theorem}

\begin{proof}
Let $M$ be $(n{+}1)$-feasible and non-$(n{+}1)$-realizable. If $M$ is forbidden,
$\ord(M)=1$. Otherwise $\ord(M)=1+\max$ over children, and every child is an
$n$-feasible non-$n$-realizable sequence, of order $\le\maxord(n)$.
\end{proof}

\begin{theorem}\label{thm:upper}
$\maxord(n)\le n-5$ for all $n\ge8$; also $\maxord(7)=3$ and $\maxord(6)=2$.
\end{theorem}

\begin{proof}
The census of Table~\ref{tab:census} is exact and exhaustive at $n=8$:
$\maxord(8)=3$. Induct with Theorem~\ref{thm:step}.
\end{proof}

The entire content of the sharp form ``$\maxord(n)=n-5$'' is therefore the lower
bound: the claim that the $+1$ step of Theorem~\ref{thm:step} is \emph{achieved} at
every level from $8$ on.

\subsection{The floor-and-towers family: every order is realized}

\begin{theorem}[Floor and towers]\label{thm:Wa}
For every $a\ge3$, the sequence
\[
W_a=\bigl(P,\;Q,\;a,\,a-1,\,\dots,\,2,\,1\bigr),\qquad
P=\tfrac{(a+1)(3a-2)}{2},\quad Q=2a+2,
\]
is $(2a{+}1)$-feasible, not $(2a{+}1)$-realizable, and has order exactly $a$.
\end{theorem}

\begin{proof}
Write $n=2a+1$, call the entries $(a,a-1,\dots,1)$ the \emph{floor}, and let
$L=\{a+1,\dots,2a+1\}$. The hypothesis $a\ge3$ enters exactly twice: it gives
$P>Q$, so that $W_a$ is sorted with pairwise distinct entries (at $a=2$ the family
degenerates to $P=Q=6$, i.e.\ the forbidden $(6,6,2,1)$), and it excludes the
one-term case $a+3=2a+1$ in the analysis of $p_j\neq Q$ below.

\emph{Feasibility.} $P+Q=\frac{(a+1)(3a-2)}{2}+2(a+1)=\frac{(a+1)(3a+2)}{2}
=\mass(L)$, so $\sum W_a=\mass(L)+\binom{a+1}{2}=\binom{2a+2}{2}$.

\emph{Floor rigidity.} In any realization, relaxed or exact, the floor blocks lie in
$\intv a$ (Lemma~\ref{lem:selfbound}) and their total demand $\binom{a+1}{2}$ equals
$\mass(\intv a)$; disjointness forces them to consume $\intv a$ exactly. Hence
tower blocks avoid $\intv a$, and in an $n$-realization the two towers must
partition $L$.

\emph{Not $n$-realizable.} $Q=2a+2$ is not a subset sum of $L$: one element is at
most $2a+1<Q$, and two or more sum to at least $(a+1)+(a+2)=2a+3>Q$.

\emph{Realizable.} $\{P\},\{Q\},\{a\},\dots,\{1\}$ are pairwise disjoint singletons
($P\neq Q$ for $a\ge3$; both exceed $a$).

\emph{Order.} Consider any node of $T(W_a,n)$ with pending element $\ell$; it has
the form (two tower values $u,v$; floor), where $u+v=\mass(\{a+1,\dots,\ell\})$ is
preserved along any path (each assignment moves in lock-step with the shrinking
$L$). Two facts:
\begin{itemize}[itemsep=1pt]
\item \emph{Any tower value $t\in\intv a$ is fatal}: the entries $\le a$ then have
total demand $\binom{a+1}{2}+t>\binom{a+1}{2}\ge\binom{s+1}{2}$ for their largest
entry $s\le a$ --- the tail volume condition~\eqref{eq:volume} fails and the node is
forbidden.
\item \emph{A $Q$-move dies immediately}: at pending $\ell\in[a+2,\,2a+1]$,
assigning $\ell$ to the tower of value $Q=2a+2$ leaves $Q-\ell\in\intv a$.
\end{itemize}
Hence while $\ell\ge a+2$, every surviving path assigns pending elements to the
$P$-tower only. After $j$ such moves the $P$-value is
$p_j=P-\sum_{i=2a+2-j}^{2a+1}i$; since $P+Q=\mass(L)$,
\[
p_{a-1}\;=\;P-\sum_{i=a+3}^{2a+1}i\;=\;P-\bigl(\mass(L)-(a+1)-(a+2)\bigr)
\;=\;(2a+3)-Q\;=\;1,
\]
and consequently $p_j=1+\sum_{i=a+3}^{2a+1-j}i$, so $p_{a-2}=a+4$. For $j\le a-2$: $p_j\ge a+4>a$;
$p_j\neq Q$, since $p_j=2a+2$ would force $\sum_{i=a+3}^{2a+1-j}i=2a+1$, and that
sum is either $0$ (empty range), or $a+3\neq2a+1$ (one term, using $a\ge3$), or at
least $(a+3)+(a+4)=2a+7$ (two or more terms); and $p_j>2a+1-j=$ pending, so the
next $P$-move is legal. Each
such node has pairwise distinct entries ($p_j$, $Q$, floor), hence is realizable by
singletons (Corollary~\ref{cor:distinct}). So the alive nodes are exactly the
$P$-path prefixes at levels $1,\dots,a-1$. From the deepest one (level $a-1$,
$P$-value $p_{a-2}=a+4$, pending $a+3$) both moves die at level $a$: the $P$-move
leaves $p_{a-1}=1\in\intv a$, the $Q$-move leaves $2a+2-(a+3)=a-1\in\intv a$. All
other branches died at levels $\le a$. Therefore $\ord(W_a)=a$.
\end{proof}

The family reproduces known extremal witnesses: $W_3=(14,8,3,2,1)$ attains
$\maxord(7)=3$ and $W_4=(25,10,4,3,2,1)$ attains $\maxord(9)=4$. Note that the
entries of $W_a$ are pairwise distinct, so the family also shows that
distinct-entry witnesses have unbounded order --- though, by
Remark~\ref{rem:incomparable}, ones transparent to the counting oracle; the
counting-blind analogue is Theorem~\ref{thm:oddtowers}. The same
construction one step earlier covers the even levels:

\begin{proposition}[Even levels]\label{prop:Vr}
For every $r\ge4$ the sequence
\[
V_r=\bigl(P',\;Q',\;r,\,r-1,\,\dots,\,2,\,1\bigr),\qquad
P'=\tfrac{r(3r+1)}{2}-(2r+2),\quad Q'=2r+2,
\]
is $2r$-feasible, not $2r$-realizable, and has order exactly $r-1$. (The hypothesis
$r\ge4$ is needed already for sortedness: $V_3=(7,8,3,2,1)$ has $P'<Q'$.)
\end{proposition}

\begin{proof}
The proof of Theorem~\ref{thm:Wa} transfers step for step; we record only the
constants and the two uses of the hypothesis. Write $n=2r$ and
$L=\{r+1,\dots,2r\}$, so $\mass(L)=\tfrac{r(3r+1)}{2}=P'+Q'$ and $V_r$ is
$n$-feasible. The hypothesis $r\ge4$ enters first through sortedness: $P'>Q'$
iff $3r^2-7r-8>0$ iff $r\ge4$, and then the entries are pairwise distinct
($P'>Q'>r$), so $V_r$ is realizable by singletons. Floor rigidity and the
non-representability of $Q'=2r+2$ as a subset sum of $L$ hold verbatim, so
$V_r$ is not $n$-realizable.

In the tree, tower values in $\intv r$ are fatal and a $Q'$-move at pending
$\ell\in[r+2,\,2r]$ leaves $Q'-\ell\in[2,r]$, so every surviving path feeds the
$P'$-tower; from $P'+Q'=\mass(L)$ the $P'$-values satisfy
$p_{r-2}=(2r+3)-Q'=1$, hence $p_j=1+\sum_{i=r+3}^{2r-j}i$ and $p_{r-3}=r+4$.
The alive nodes are the $P'$-path prefixes at levels $1,\dots,r-2$, exactly as
before; the check $p_j\neq Q'$ needs $r+3\neq2r+1$ in its one-term case --- the
second use of $r\ge4$. From the deepest alive node (level $r-2$, pending
$r+3$) both moves die at level $r-1$: $p_{r-2}=1$ and $Q'-(r+3)=r-1$ lie in
$\intv r$. Therefore $\ord(V_r)=r-1$.
\end{proof}

\begin{corollary}\label{cor:unbounded}
$\maxord(n)\ge\lfloor(n-1)/2\rfloor$ for all $n\ge7$; in particular
$\maxord(n)\to\infty$. Moreover every positive integer is the order of some
sequence: $(1,1,1)$ at $n=2$ has order $1$; $(8,8,3,2)$ at $n=6$ has order $2$;
$W_a$ has order $a$ for $a\ge3$.
\end{corollary}

\begin{proof}
For odd $n=2a+1\ge7$, Theorem~\ref{thm:Wa} gives $\ord(W_a)=a=(n-1)/2$; for even
$n=2r\ge8$, Proposition~\ref{prop:Vr} gives $\ord(V_r)=r-1=\lfloor(n-1)/2\rfloor$.
\end{proof}

The gap between $\lfloor(n-1)/2\rfloor$ and $n-5$ is real: from $n=10$ on the true
maximal witnesses are not floor-and-towers sequences --- their obstructions couple
mid-sized repeated entries rather than a rigid floor, which is exactly what resists
closed-form analysis (Section~\ref{subsec:horizon}).

\subsection{Lifting: $\maxord(n)=n-5$ for $8\le n\le31$}

\begin{proposition}[Lift exactness]\label{prop:lift}
Suppose the set $S_d$ of order-$d$ witnesses at level $n$ is known exactly and
$\maxord(n)=d$. Then the order-$(d{+}1)$ witnesses at level $n+1$ are exactly the
sequences obtained from some $W\in S_d$ by adding $n+1$ to one entry or appending
$n+1$ as a new entry, that are simultaneously (i) realizable (relaxed sense) and
(ii) not $(n{+}1)$-realizable.
\end{proposition}

\begin{proof}
Let $M$ be an order-$(d{+}1)$ witness at level $n+1$. It is realizable (its
order exceeds $1$), and $\ord(M)=1+\max_{C}\ord(C)$ forces some child $C$ of
order $d$, hence $C\in S_d$ since $\maxord(n)=d$. That child arose by assigning
the pending element $n+1$ to an entry of $M$, so $M$ is recovered from $C$ by
the inverse move: adding $n+1$ back onto one entry, or appending $n+1$ as a new
entry when the assignment consumed its entry exactly. Conversely, a sequence
$M$ obtained this way that satisfies (i) and (ii) has a child in $S_d$, so
$\ord(M)\ge d+1$; and $\ord(M)\le\maxord(n+1)\le d+1$ by
Theorem~\ref{thm:step}.
\end{proof}

Starting from the census fact that $(24,16,11,10,10,4,2,1)$ is the unique order-$7$
witness at $n=12$, the scan gives:

\begin{table}[ht]
\centering
\caption{Maximal-order witnesses by exact lifting (\texttt{order\_hierarchy.py}); the scan is
exhaustive, so the counts are exact, for all $n\le23$.}
\label{tab:lift}
\begin{tabular}{@{}lrrrrrrrrrrr@{}}
\toprule
$n$ & 13 & 14 & 15 & 16 & 17 & 18 & 19 & 20 & 21 & 22 & 23 \\
\midrule
order & 8 & 9 & 10 & 11 & 12 & 13 & 14 & 15 & 16 & 17 & 18 \\
count & 3 & 18 & 48 & 24 & 57 & 72 & 153 & 845 & 2194 & 6845 & 26453 \\
\bottomrule
\end{tabular}
\end{table}

\noindent Hence $\maxord(n)=n-5$ holds, with the witness sets known exhaustively,
for all $8\le n\le23$. The three order-$8$ witnesses at $n=13$ display the three
lift mechanisms: $(24,24,16,10,10,4,2,1)$ (doubling the maximum),
$(37,16,11,10,10,4,2,1)$ (growing the maximum), and $(24,16,13,11,10,10,4,2,1)$
(appending the new element). At $n=23$ every witness has $m_1\ge n+2$ and
$m_2\ge n+1$, so Theorems~\ref{thm:swap} and~\ref{thm:secondentry} stay sharp.

Beyond $n=23$ exhaustive enumeration becomes infeasible, but the certification
implicit in the proof of Proposition~\ref{prop:lift} is \emph{per witness} and does
not require the level to be exhausted: once $\maxord(n)=n-5$ is established, any
un-assignment $M$ of a known order-$(n{-}5)$ witness that is realizable and not
$(n{+}1)$-realizable satisfies $\ord(M)\ge n-4$ (via that child) and
$\ord(M)\le\maxord(n+1)\le n-4$ (Theorem~\ref{thm:step}), so $\ord(M)=n-4$ and
$\maxord(n+1)=n-4$ exactly; only the witness \emph{counts} are lost. A
capped-frontier scan retaining at most $150$ witnesses per level (\texttt{order\_hierarchy.py};
total running time about $23$ minutes) continues the chain through $n=31$ (order
$26$), e.g.\ $(272,33,25,24,24,24,24,16,16,15,10,10,2,1)$. Hence
\[
\maxord(n)=n-5\qquad\text{for all } 8\le n\le31.
\]

\subsection{Horizon theorems: fixed cores cannot finish the job}\label{subsec:horizon}

The natural attack on the lower bound is a family $M_n$ built from a fixed finite
core $F$ by absorbing each new element. The first absorption mode provably dies;
the second dies in every instance we have checked.

\begin{theorem}[Top-singleton reduction]\label{thm:topsingleton}
If an $n$-feasible $M$ contains an entry equal to $n$, then $M$ is $n$-realizable
iff $M$ minus one entry $n$ is $(n{-}1)$-realizable.
\end{theorem}

\begin{proof}
($\Leftarrow$) Add the block $\{n\}$. ($\Rightarrow$) In a partition of $\intv n$
the element $n$ lies in some block $B_e$ realizing an entry $e$. If $e=n$, delete
both. Otherwise $e>n$; let $B$ be the block of an entry $n$ (so $n\notin B$) and
exchange: $B_e\leftarrow(B_e\setminus\{n\})\cup B$ preserves sums and disjointness
and gives the entry $n$ the block $\{n\}$; delete both.
\end{proof}

\begin{corollary}[Interval extensions]\label{cor:interval}
For any $c$-feasible non-$c$-realizable $F$ and every $n\ge c$, the sequence
$F\cup\{c+1,c+2,\dots,n\}$ is $n$-feasible and not $n$-realizable.
\end{corollary}

This looks like a royal road: it manufactures non-realizable sequences at every
level from one finite seed, and if every $M_m=F\cup\{c+1..m\}$ were realizable in
the relaxed sense, the descending chain would give
$\ord(M_n)=(n-c)+\ord(F)$. But:

\begin{theorem}[Horizon]\label{thm:horizon}
Let $F$ be $c$-feasible and non-$c$-realizable, and $M_m=F\cup\{c+1,\dots,m\}$.
Then $M_m$ is forbidden (relaxed sense) for every $m\ge\max(F)-1$. Hence the
interval extension certifies order $m-5$ at most up to $m=\max(F)-2$.
\end{theorem}

\begin{proof}
For $m\ge\max(F)$ the relaxed ground set is $\intv m$ and the entry mass is the full
$\binom{m+1}{2}$ (feasibility): disjoint blocks carrying the whole mass form a
partition, so relaxed realizability degenerates to $m$-realizability --- false by
Corollary~\ref{cor:interval}. (This is the boundary collapse of
Theorem~\ref{thm:boundary}.) For $m=\max(F)-1$ the ground set $\intv{m_1}$ contains
exactly one element exceeding $m$, so a relaxed realization would use at most one
overflow element, and Theorem~\ref{thm:swap} would make $M_m$
$m$-realizable --- false again.
\end{proof}

As a demonstration, $F=(24,24,16,16,16,16,10,10,4)$ --- a maximal witness at
$n=16$, one of the $24$ in Table~\ref{tab:lift} --- gives $G_m=F\cup\{17..m\}$ of
order $m-5$ for $m=16,\dots,22$, six levels for free, and $G_m$ forbidden from
$m=23=\max(F)-1$ on, exactly as predicted. We expect, but have not proved, that the
second absorption mode --- growing two towers over a fixed tail $C$ --- dies for the
same underlying reason: heuristically, non-realizable tower splits should exist
only while $n<\max(C)$, since once the ground set reaches the largest tail entry
the exchange of Theorem~\ref{thm:topsingleton} frees it. We record this as a
verified instance rather than a theorem: for the tail
$(24,18,16,16,15,10,10,4,2,1)$, non-realizable splits exist at $n=22,23$ and none
from $n=24$ on.

Together with $m_1\ge n+2$, $m_2\ge n+1$, the moral is that a maximal-order witness
cannot carry any fixed finite structure: its top entries must grow linearly and its
``tail'' must continually re-form just below the pending element. This is visible
in the data --- the witness shapes metamorphose at $n=23$ and again after
$n=24$ --- and is why we expect the following to require a genuinely scale-invariant
construction.

\begin{conjecture}\label{conj:lower}
$\maxord(n)=n-5$ for all $n\ge8$.
\end{conjecture}

\section{Complexity}\label{sec:complexity}

Neither decision problem in this paper has known complexity status.

\textbf{$n$-realizability.} No NP-hardness proof and no polynomial algorithm is
known, and the extremal literature does not pose the question. Two subtleties
deserve note. First, the natural input is $k$ numbers ($O(k\log m_1)$ bits) while a
witness partition of $\intv n$ has size $\Theta(n)$, which can be exponential in the
input size; even NP-membership is therefore not automatic in the succinct encoding,
though it is immediate in unary. Second, the standard sources of strong NP-hardness
for prescribed-sums partitioning --- 3-Partition~\cite{GareyJohnson79} and multiple
subset sums~\cite{CapraraKP00,CieliebakEPS08} --- rely on engineered multisets of
items; no known reduction forces the item set to be exactly $\intv n$, each element
once. For fixed $k$ the unary problem is polynomial via $k$-Subset-Sum dynamic
programming~\cite{AntonopoulosPPV23}; the announced characterization
of~\cite{LladoMoragas09} decides it for $n\ge(4k)^3$; the open regime is $k$
growing with $n$.

\textbf{Relaxed realizability.} By Corollary~\ref{cor:core} the problem reduces to
its core, and under Conjecture~\ref{conj:pairs} to the full rainbow matching
Question~\ref{q:rainbow}. While the worst case remains open, the structure theory
of Sections~\ref{sec:heredity}--\ref{sec:minimal} yields two parameterized upper
bounds with no analogue for $n$-realizability. The first makes precise the sense
in which the relaxation is computationally lighter: the singleton rule collapses
the exponent from the length $k$ to the number of excess copies $r=k-d$.

\begin{proposition}[Parameterized core dynamic programming]\label{prop:coredp}
Let $M$ have $k$ entries and $d$ distinct values, and let $v_1\ge\dots\ge v_r$,
$r=k-d$, be the values of the excess copies listed with multiplicity (so $v_1$ is
the largest repeated value). Realizability of $M$ can be decided deterministically
in time
\[
O\Bigl(k+v_1\,r\prod_{j=1}^{r}(v_j+1)\Bigr)\;\subseteq\;
O\bigl(k+r\,v_1^{\,r+1}\bigr),
\]
and in randomized time $\Otilde(k+v_1^{\,r})$. In particular sequences with
pairwise distinct entries are decided in time $O(k)$, and the exponent is
governed by repetition alone, not by length.
\end{proposition}

\begin{proof}
By the First-Repeat Reduction (Theorem~\ref{thm:firstrepeat}) we may delete every
entry above the first repeated value; this changes neither $r$ nor the excess
values, and the surviving sequence $T$ has maximum entry $v_1$. By
Corollary~\ref{cor:core}, $T$ is realizable iff there are pairwise disjoint
subsets of $A=\intv{v_1}\setminus V(T)$ with sums $v_1,\dots,v_r$, where $V(T)$
is the set of distinct values of $T$; the core's size-$\ge2$ constraint is
automatic, since $v_j\notin A$ forces any subset of $A$ summing to $v_j$ to have
at least two elements. This is an instance of $r$-Subset-Sum on $|A|<v_1$ items.
The Bellman dynamic program over state vectors $(s_1,\dots,s_r)$ with
$0\le s_j\le v_j$ --- each element of $A$ in turn is either discarded or added to
one unfinished block --- runs in time $O(|A|\cdot r\prod_j(v_j+1))$; sorting,
the volume check, and stripping cost $O(k+v_1)$. The randomized bound is the
$k$-Subset-Sum algorithm of~\cite{AntonopoulosPPV23} (time $\Otilde(n+t^k)$ on
$n$ items with $k$ targets bounded by $t$) applied to the same instance.
\end{proof}

For comparison, the same machinery applied to $n$-realizability gives
$\Otilde(n+m_1^{\,k})$: run $k$-Subset-Sum on the ground set $\intv n$ with all
$k$ targets --- disjoint subsets meeting all targets automatically cover $\intv n$,
since by $n$-feasibility their sums exhaust $\mass(\intv n)$ --- and no collapse
of the exponent is available, because the exchange proving
Theorem~\ref{thm:singleton} moves elements \emph{out of use}, exactly the freedom
the covering constraint removes. The gap --- exponent $r$ versus $k$, ground set
$\intv{v_1}$ versus $\intv n$ --- quantifies the relaxation's lightness. We also
note $2r\le|A|<v_1$ whenever the count condition of Section~\ref{subsec:hall}
holds, so instances reaching the dynamic program always have
$r\le\min(k-1,\lfloor v_1/2\rfloor)$.

The second bound is conditional on the pairs conjecture and covers the opposite
regime: few distinct values, arbitrary repetition.

\begin{proposition}[Reduction to Exact Matching]\label{prop:exactmatching}
Assume Conjecture~\ref{conj:pairs}. Then realizability is decidable by a
randomized algorithm with one-sided error in time
$\mathrm{poly}(m_1)\cdot(k+1)^{O(d)}$, where $d$ is the number of distinct
values; in particular, in randomized polynomial time when $d$ is bounded.
Acceptance is sound unconditionally: the algorithm accepts only by exhibiting a
realization with all $|X_i|\le2$.
\end{proposition}

\begin{proof}
Under Conjecture~\ref{conj:pairs}, $M$ is realizable iff one can select $\mu(v)$
pairwise disjoint members of the menu $E_v$ for every $v\in V$
(Section~\ref{subsec:rainbow}); recall from the discussion there that the
singleton options must remain selectable. A pairs realization uses
$2k-s\le m_1$ elements, where $s\le d$ is its number of singleton blocks, so if
$2k-d>m_1$ we may reject outright; assume $2k-d\le m_1$. Build a graph $G$ on the
vertex set $\intv{m_1}\cup\{s_v:v\in V\}\cup D$: for each $v\in V$ the edges
$\{x,v-x\}$, $1\le x<v/2$, together with the auxiliary edge $\{v,s_v\}$
representing the singleton $\{v\}$; and a set $D$ of $m_1+d-2k$ dummy vertices
adjacent to every non-dummy vertex. Note every non-dummy edge has a well-defined
\emph{color}, the common sum $v$ of its endpoints (with $\{v,s_v\}$ colored $v$),
and $|V(G)|=2(m_1+d-k)$ is even. Selections of $\mu(v)$ pairwise disjoint members
of each $E_v$ correspond exactly to perfect matchings of $G$ containing $\mu(v)$
edges of color $v$ for every $v$: a selection is a matching of $k$ colored edges
covering $2k$ non-dummy vertices, and the $m_1+d-2k$ uncovered non-dummies pair
off with the dummies (the dummy-to-non-dummy graph is complete bipartite);
conversely the dummies of a perfect matching absorb all but $2k$ non-dummy
vertices, which are perfectly matched by $k$ colored edges. Now index
$V=\{u_1,\dots,u_d\}$ and give edges of color $u_j$ the weight $(k+1)^{j-1}$,
dummy edges the weight $0$. A perfect matching contains $k$ colored edges, so its
color counts $c_j\le k$ are digits base $k+1$: its total weight equals
$W=\sum_j\mu(u_j)(k+1)^{j-1}$ iff $c_j=\mu(u_j)$ for all $j$. Deciding whether
$G$ has a perfect matching of total weight exactly $W$ --- and extracting one ---
is the Exact Matching problem~\cite{PapadimitriouYannakakis82} with weights
bounded by $(k+1)^{d}$, solvable in randomized time polynomial in $|V(G)|$ and
the weight bound by the Pfaffian--isolation technique of Mulmuley, Vazirani and
Vazirani~\cite{MulmuleyVaziraniVazirani87}: substitute $y^{w(e)}$ times random
scalars for the edge variables of the Tutte matrix and recover the coefficients
of the determinant as a polynomial in $y$ by interpolation. If a matching of
weight $W$ is found, its colored edges are a pairs realization --- a witness valid
without any conjecture; if none exists, Conjecture~\ref{conj:pairs} certifies $M$
forbidden.
\end{proof}

Propositions~\ref{prop:coredp} and~\ref{prop:exactmatching} parameterize
complementary regimes: the first is efficient when repeats are few ($r=k-d$
small), the second when the value spectrum is narrow ($d$ small). The regime
where both parameters grow --- many distinct repeated values --- is exactly where
the minimal forbidden catalog of Section~\ref{sec:minimal} and the tower families
of Section~\ref{sec:order} live, and is the content of
Question~\ref{q:rainbow}. Conjecture~\ref{conj:pairs} would also settle succinct
NP-membership for the relaxed problem: a pairs realization is a certificate of
$O(k\log m_1)$ bits, polynomial in the binary encoding --- in contrast with
$n$-realizability above, where no succinct certificate is known. Beyond these
regimes, the hardness landscape for unstructured full rainbow
matchings (Section~\ref{subsec:rainbow}) shows that ``all color classes are
matchings'' cannot by itself give tractability; conversely, no hardness reduction
we are aware of can produce the interval menus $E_v$. We also note that Exact
Matching itself~\cite{PapadimitriouYannakakis82}, in randomized NC yet
not known to be in P, illustrates that prescribed-count-per-class matching
problems can occupy this intermediate zone.

In practice the relaxed problem is light: the layered solver of
Appendix~\ref{app:solver} classifies all $7.4$ million sequences with $m_1\le12$
in seconds, with the volume conditions~\eqref{eq:volume} alone settling $97\%$ of
instances. The lightness is not uniform, however: randomized
stress tests (Appendix~\ref{app:solver}) locate near-tight instances with many
distinct repeated values on which both implementations stall, concentrating
empirical hardness in the same two-parameter regime identified above. Whether
the practical lightness of typical instances reflects polynomial-time
decidability is the central algorithmic question left open.

\section{Open problems}\label{sec:open}

\begin{enumerate}[itemsep=2pt]
\item \textbf{Pairs.} Prove Conjecture~\ref{conj:pairs}: every realizable sequence
has a realization with all $|X_i|\le2$. (Verified for $m_1\le18$; for length
$k\le5$ up to $m_1\le120$ and $k\le6$ up to $m_1\le40$; known for distinct-entry,
constant, and length-$\le3$ sequences by Proposition~\ref{prop:pairsspecial}.
The remaining difficulty is reduced to orienting the pair swap of
Proposition~\ref{prop:saturation}.)
\item \textbf{Structured rainbow matching / complexity.} Resolve
Question~\ref{q:rainbow}; more broadly, determine the complexity of relaxed
realizability and of $n$-realizability with input in unary.
(Propositions~\ref{prop:coredp} and~\ref{prop:exactmatching} settle the regimes
$r=O(1)$ and --- conditionally --- $d=O(1)$; the open case is $r,d$ both
growing.)
\item \textbf{Frontier.} Prove Conjecture~\ref{conj:frontier}:
$n_{\max}(k)=4(k-2)$ with unique extremal witness $F_{k-2}$ for $k\ge4$.
\item \textbf{Sharp lower bound.} Prove Conjecture~\ref{conj:lower}:
$\maxord(n)=n-5$ for all $n\ge8$ (known for $8\le n\le31$).
By Theorems~\ref{thm:topsingleton}
and~\ref{thm:horizon} any proof requires a family whose cores grow with $n$; the
census suggests towers of size $\Theta(n)$ over a geometric ladder of repeated
mid-sized entries.
\item \textbf{Order two.} Characterize the order-$2$ sequences. By
Theorem~\ref{thm:swap} all their realizations need at least two overflow elements,
and the failure is a coverage/subset-sum obstruction in $\intv n$ after forced
small blocks. The family $D_n$ (Proposition~\ref{prop:Dn}) supplies distinct-entry
order-$2$ witnesses at every level $n\ge6$; a characterization should in particular
explain which order-$2$ sequences the counting conditions of
Proposition~\ref{prop:counting} detect.
\item \textbf{Distinct entries and counting.} The odd-towers family $O_n$
satisfies every counting condition of Proposition~\ref{prop:counting} yet has
$\ord(O_n)=n-8$ for even $24\le n\le32$ (Theorem~\ref{thm:oddtowers}); is this
exact for all even $n\ge24$? Distinct-entry sequences attain $\maxord(n)$ for
$n\le9$ but fall exactly one short for $10\le n\le12$
(Section~\ref{subsec:census}): can distinct-entry --- or counting-blind ---
sequences attain $\maxord(n)$ infinitely often?
\item \textbf{Certificates.} Is there a polynomial certificate of
non-$n$-realizability, or is the unbounded order hierarchy
(Corollary~\ref{cor:unbounded}) evidence of genuine hardness beyond the
relaxation?
\item \textbf{Enumeration.} Determine the growth rate of the number of minimal
forbidden sequences with maximum entry $n$ (Table~\ref{tab:catalog}).
\end{enumerate}

\appendix

\section{The exact solver}\label{app:solver}

The layered algorithm \texttt{realizable\_fast} (\texttt{layered\_solver.py}) is exact
regardless of Conjecture~\ref{conj:pairs}: (1) tail volume check~\eqref{eq:volume},
$O(k)$, rejecting $97\%$ of all sequences outright; (2) singleton stripping
(Corollary~\ref{cor:core}), provably safe; (3) two greedy passes, sound when they
find a witness; (4) complete pairs-only search (sound; conjecturally complete);
(5) complete memoized backtracking with symmetry breaking and per-node volume
pruning. It agrees with brute force on all $3{,}002$ sequences with $m_1\le8$,
$k\le6$; classifies all $7.4$M sequences with $m_1\le12$ in ${\approx}25$\,s; and
solves random instances with $m_1$ up to $100$ in well under a millisecond. The
worst case observed is adversarial repetition (ten copies of $50$ with near-tight
small entries): $3.4$\,s. The dynamic program of Proposition~\ref{prop:coredp} is
implemented separately in \texttt{core\_dp.py}, whose self-test re-verifies
agreement with \texttt{realizable\_fast} on every sequence with $m_1\le9$
($246{,}232$ sequences, including all volume-infeasible ones up to a margin).

The two solvers are complementary in guarantees rather than in practice
(\texttt{bench\_solvers.py}). Where the dynamic program is tractable --- few
excess copies --- the greedy layers already succeed in microseconds and outrun
it: random instances with $r=2$ and $60\le m_1\le100$ take under a millisecond
for the layered solver against several milliseconds for the dynamic program,
which always pays for its full table. Conversely, the empirically hard
instances for the backtracking layer are near-tight sequences with many excess
copies: in a randomized stress test over volume-feasible instances with
$m_1\le60$ and multiplicities up to $4$, every instance on which
\texttt{realizable\_fast} exceeded $50$\,ms --- some exceeding a $2$\,s guard ---
had $r\ge5$, typically $11$--$17$, where the state space $r\prod_j(v_j+1)$ of
Proposition~\ref{prop:coredp} ranges from $10^{12}$ to $10^{36}$. The two
methods thus share a single hard regime --- many distinct repeated values ---
precisely the regime left open after Propositions~\ref{prop:coredp}
and~\ref{prop:exactmatching}.

\section{Census of the order hierarchy, $n\le12$}\label{app:census}

\begin{table}[ht]
\centering
\caption{Classification of all $n$-feasible sequences, $n\le12$
(\texttt{assignment\_tree.py}).}
\label{tab:census}
\footnotesize
\begin{tabular}{@{}rrrrrlr@{}}
\toprule
$n$ & feasible & realizable & $n$-realizable & order $\ge2$ & orders (count) & $\maxord(n)$ \\
\midrule
$\le5$ & 233 & 57 & 57 & 0 & --- & 1 \\
6 & 792 & 112 & 109 & 3 & 2:\,3 & 2 \\
7 & 3{,}718 & 355 & 340 & 15 & 2:\,13,\ 3:\,2 & 3 \\
8 & 17{,}977 & 1{,}190 & 1{,}116 & 74 & 2:\,65,\ 3:\,9 & 3 \\
9 & 89{,}134 & 4{,}091 & 3{,}744 & 347 & 2:\,245,\ 3:\,93,\ 4:\,9 & 4 \\
10 & 451{,}276 & 14{,}461 & 12{,}981 & 1{,}480 & 2:\,937,\ 3:\,427,\ 4:\,115,\ 5:\,1 & 5 \\
11 & 2{,}323{,}520 & 51{,}965 & 45{,}722 & 6{,}243 & 2:\,3237,\ 3:\,2201,\ 4:\,672,\ 5:\,130,\ 6:\,3 & 6 \\
12 & 12{,}132{,}164 & 190{,}295 & 165{,}247 & 25{,}048 & 2:\,11622,\ 3:\,8276,\ 4:\,4003,\ 5:\,1054,\ 6:\,92,\ 7:\,1 & 7 \\
\bottomrule
\end{tabular}
\end{table}

For $n\le5$ every feasible non-$n$-realizable sequence is already forbidden, so
$\maxord(n)=1$ there; the ``orders (count)'' column lists only the sequences of
order $\ge2$.

\begin{thebibliography}{99}

\bibitem{AlaviBCEO87}
Y.~Alavi, A.~J. Boals, G.~Chartrand, P.~Erd\H{o}s, O.~R. Oellermann,
The ascending subgraph decomposition problem,
\emph{Congr. Numer.} \textbf{58} (1987), 7--14.

\bibitem{AndoGervacioKano90}
K.~Ando, S.~Gervacio, M.~Kano,
Disjoint subsets of integers having a constant sum,
\emph{Discrete Math.} \textbf{82} (1990), 7--11.

\bibitem{AntonopoulosPPV23}
A.~Antonopoulos, A.~Pagourtzis, S.~Petsalakis, M.~Vasilakis,
Faster algorithms for $k$-Subset Sum and variations,
\emph{J. Comb. Optim.} \textbf{45} (2023), art. 24.

\bibitem{ArocaLlado18}
J.~M. Aroca, A.~Lladó,
On star forest ascending subgraph decomposition,
\emph{Electron. J. Combin.} \textbf{25}(1) (2018), \#P1.67. arXiv:1512.02161.

\bibitem{BuchelGW18}
A.~Büchel, U.~Gilleßen, K.-U. Witt,
An output-sensitive algorithm to partition a sequence of integers into subsets with
equal sums,
\emph{Discrete Math. Theor. Comput. Sci.} \textbf{20}(2) (2018).

\bibitem{CapraraKP00}
A.~Caprara, H.~Kellerer, U.~Pferschy,
The multiple subset sum problem,
\emph{SIAM J. Optim.} \textbf{11} (2000), 308--319.

\bibitem{ChenFuWangZhou05}
F.-L. Chen, H.-L. Fu, Y.~Wang, J.~Zhou,
Partition of a set of integers into subsets with prescribed sums,
\emph{Taiwanese J. Math.} \textbf{9} (2005), 629--638.

\bibitem{CichaczGorlich18}
S.~Cichacz, A.~Görlich,
Constant sum partition of sets of integers and distance magic graphs,
\emph{Discuss. Math. Graph Theory} \textbf{38} (2018), 97--106.

\bibitem{CieliebakEPS08}
M.~Cieliebak, S.~Eidenbenz, A.~Pagourtzis, K.~Schlude,
On the complexity of variations of Equal Sum Subsets,
\emph{Nordic J. Comput.} \textbf{14} (2008), 151--172.

\bibitem{Davies59}
R.~O. Davies,
On Langford's problem (II),
\emph{Math. Gaz.} \textbf{43} (1959), 253--255.

\bibitem{deFigueiredoMSS22}
C.~M.~H. de~Figueiredo, A.~A. de~Melo, D.~Sasaki, A.~Silva,
Revising Johnson's table for the 21st century,
\emph{Discrete Appl. Math.} \textbf{323} (2022), 184--200.

\bibitem{EnomotoKano95}
H.~Enomoto, M.~Kano,
Disjoint odd integer subsets having a constant even sum,
\emph{Discrete Math.} \textbf{137} (1995), 189--193.

\bibitem{FuHu94}
H.-L. Fu, W.-H. Hu,
Disjoint odd integer subsets having a constant odd sum,
\emph{Discrete Math.} \textbf{128} (1994), 143--150.

\bibitem{GareyJohnson79}
M.~R. Garey, D.~S. Johnson,
\emph{Computers and Intractability: A Guide to the Theory of NP-Completeness},
W.~H. Freeman, 1979.

\bibitem{HommelsheimJM26}
F.~Hommelsheim, P.~Jehmlich, M.~Mühlenthaler,
A complexity dichotomy for generalized rainbow matchings based on color classes,
arXiv:2604.21025 (2026).

\bibitem{HooryKotlar26}
S.~Hoory, D.~Kotlar,
Solvable and unsolvable instances of the equal sum partition problem,
arXiv:2605.06071 (2026).

\bibitem{ItaiRodehTanimoto78}
A.~Itai, M.~Rodeh, S.~L. Tanimoto,
Some matching problems for bipartite graphs,
\emph{J. ACM} \textbf{25} (1978), 517--525.

\bibitem{KLPS25}
K.~Katsamaktsis, S.~Letzter, A.~Pokrovskiy, B.~Sudakov,
Ascending subgraph decomposition,
\emph{J. Combin. Theory Ser. B} \textbf{173} (2025), 14--44. arXiv:2308.11613.

\bibitem{KelkStamoulis19}
S.~Kelk, G.~Stamoulis,
Integrality gaps for colorful matchings,
\emph{Discrete Optim.} \textbf{32} (2019), 73--92.

\bibitem{Langford58}
C.~D. Langford,
Problem,
\emph{Math. Gaz.} \textbf{42} (1958), 228.

\bibitem{LePfender14}
V.~B. Le, F.~Pfender,
Complexity results for rainbow matchings,
\emph{Theoret. Comput. Sci.} \textbf{524} (2014), 27--33.

\bibitem{LladoMoragas09}
A.~Lladó, J.~Moragas,
On the sumset partition problem,
\emph{Electron. Notes Discrete Math.} \textbf{34} (2009), 15--19.

\bibitem{LladoMoragas12}
A.~Lladó, J.~Moragas,
On the modular sumset partition problem,
\emph{European J. Combin.} \textbf{33} (2012), 427--434.

\bibitem{MaZhouZhou94}
K.~Ma, H.~Zhou, J.~Zhou,
On the ascending star subgraph decomposition of star forests,
\emph{Combinatorica} \textbf{14} (1994), 307--320.

\bibitem{Moragas10}
J.~Moragas,
\emph{Graph labelings and decompositions by partitioning sets of integers},
Ph.D. thesis, Universitat Polit\`ecnica de Catalunya, Barcelona, 2010.

\bibitem{MulmuleyVaziraniVazirani87}
K.~Mulmuley, U.~V. Vazirani, V.~V. Vazirani,
Matching is as easy as matrix inversion,
\emph{Combinatorica} \textbf{7} (1987), 105--113.

\bibitem{OEIS}
OEIS Foundation Inc.,
\emph{The On-Line Encyclopedia of Integer Sequences},
\url{https://oeis.org} (accessed July 2026).

\bibitem{PapadimitriouYannakakis82}
C.~H. Papadimitriou, M.~Yannakakis,
The complexity of restricted spanning tree problems,
\emph{J. ACM} \textbf{29} (1982), 285--309.

\bibitem{Skolem57}
T.~Skolem,
On certain distributions of integers in pairs with given differences,
\emph{Math. Scand.} \textbf{5} (1957), 57--68.

\bibitem{TanTeh24}
T.~S. Tan, W.~C. Teh,
A note on graph burning of path forests,
\emph{Discrete Math. Theor. Comput. Sci.} \textbf{26}(3) (2024). arXiv:2312.10914.

\end{thebibliography}

\end{document}
